\documentclass[11pt,twocolumn,a4paper]{article}
\usepackage[utf8]{inputenc}
\usepackage[T1]{fontenc}

% Page layout
\usepackage[margin=2.2cm]{geometry}

% Essential packages
\usepackage{amsmath,amssymb,amsthm,mathtools}
\usepackage{bm}
\usepackage{graphicx}
\usepackage{hyperref}
\usepackage{natbib}
\usepackage{physics}
\usepackage{booktabs}
\usepackage{array}
\usepackage{multirow}
\usepackage{enumitem}
\usepackage{float}
\usepackage{longtable}
\usepackage{tabularx}

% Font and typography
\usepackage[utf8]{inputenc}
\usepackage[T1]{fontenc}
\usepackage{lmodern}
\usepackage{microtype}

\usepackage{caption}
\captionsetup{
  font=small,          % or footnotesize, scriptsize
  labelfont=footnotesize,        % Keep "Figure 1" bold
  textfont=it          % Optional: italicize caption text
}

% Theorem environments
\newtheorem{theorem}{Theorem}[section]
\newtheorem{lemma}[theorem]{Lemma}
\newtheorem{proposition}[theorem]{Proposition}
\newtheorem{corollary}[theorem]{Corollary}
\theoremstyle{definition}
\newtheorem{definition}[theorem]{Definition}
\newtheorem{axiom}{Axiom}
\newtheorem{example}[theorem]{Example}
\theoremstyle{remark}
\newtheorem{remark}[theorem]{Remark}

% Hyperref setup
\hypersetup{
    colorlinks=true,
    linkcolor=blue,
    citecolor=blue,
    urlcolor=blue
}

% Custom commands
\newcommand{\kB}{k_{\mathrm{B}}}
\newcommand{\Sk}{S_{\mathrm{k}}}
\newcommand{\St}{S_{\mathrm{t}}}
\newcommand{\Se}{S_{\mathrm{e}}}
\newcommand{\Sspace}{\mathcal{S}}
\newcommand{\Depth}{\mathcal{M}}
\newcommand{\Real}{\mathbb{R}}
\newcommand{\Natural}{\mathbb{N}}
\newcommand{\Integer}{\mathbb{Z}}
\newcommand{\Hilbert}{\mathcal{H}}
\newcommand{\tem}{\tau_{\mathrm{em}}}
\newcommand{\Aki}{A_{ki}}
\newcommand{\qmax}{q_{\mathrm{max}}}
\newcommand{\etamax}{\eta_{\mathrm{max}}}

\begin{document}

\title{\textbf{Harmonic Molecular Resonator: \\
Reflexive Clock Networks from Absorption-Emission \\
Intervals in Bounded Phase Space}}

\author{
Kundai Farai Sachikonye\\[0.3em]
AIMe Registry for Artificial Intelligence\\[0.3em]
\texttt{kundai.sachikonye@bitspark.com}
}

\date{\today}

\maketitle

\begin{abstract}
We derive harmonic molecular networks with circulating light from bounded phase space geometry. The absorption-emission interval $\tem$ of a molecular electronic transition provides a physical tick---a natural clock whose period is determined by partition depth, not by external timing. Virtual molecules with successively shorter emission lifetimes $\tem^{(k)} < \tem^{(k-1)}$ form a nested tick hierarchy. Each tier builds a tree structure from parent-child emission relationships. Trees become networks when harmonically close nodes across adjacent trees are joined: if $\omega_i/\omega_j = p/q$ (low-order rational ratio), the nodes share a harmonic edge. Closed loops in this harmonic network constitute virtual resonant cavities---light circulates without walls, sustained by categorical coupling rather than spatial confinement. The network is self-clocking (each circulation cycle IS a tick) and self-validating (each loop cross-checks partition coordinates through independent paths). We validate on six molecules (H$_2$, CO, H$_2$O, CO$_2$, CH$_4$, C$_6$H$_6$) using NIST spectroscopic data, demonstrating exact harmonic structure, closed-loop formation, and circulation periods consistent with known spectroscopic intervals, with zero adjustable parameters.

\medskip
\noindent\textbf{Keywords:} harmonic molecular networks, absorption-emission intervals, self-clocking systems, virtual resonant cavities, categorical coupling, Einstein A coefficient, nested tick hierarchy, Poincaré recurrence, circulating light, molecular spectroscopy

\end{abstract}

%==============================================================================
% PART I: FOUNDATIONS
%==============================================================================

\part{Foundations}

%==============================================================================
\section{Introduction}
\label{sec:introduction}
%==============================================================================

\subsection{The Ticking Problem}

The companion paper \citep{sachikonye2025a} derived the structure of chemical elements from a single axiom---the Bounded Phase Space Law---and proved that the computer is a physical spectrometer through the Triple Equivalence Theorem. A recurring objection to the categorical measurement framework is the \emph{ticking problem}: if measurement is categorical state enumeration, what provides the tick? What is the physical clock that advances the categorical state counter? Without an answer, the framework appears to require an external timekeeper---a homunculus---undermining its claim to self-sufficiency.

This paper answers the ticking problem. The answer is not an external clock but an intrinsic one: the absorption-emission interval of a molecular electronic transition IS the tick. A molecule absorbs a photon, enters an excited state, and spontaneously emits after a time $\tem = 1/\Aki$, where $\Aki$ is the Einstein A coefficient \citep{einstein1917}. This interval is determined entirely by the molecular partition structure---the transition dipole moment, the partition coordinates, and the frequency of the mode. No external timing reference is needed. The tick is built into the molecule.

\subsection{From Atoms to Molecules}

Paper~1 \citep{sachikonye2025a} derived atoms: isolated bounded systems characterized by the partition coordinate system $(n, l, m, s)$ with shell capacity $C(n) = 2n^2$. This paper derives molecules: composite bounded systems in which multiple oscillatory modes interact through harmonic proximity. The transition from atom to molecule is the transition from tree to network---from hierarchical parent-child relationships to a topology admitting closed loops.

The central result is that closed loops in the harmonic molecular network constitute \emph{virtual resonant cavities}: light circulates within the molecule without mirrors, walls, or spatial confinement. The confinement is categorical---the loop closes because the partition coordinates return to their initial values after traversing the loop. This is a fundamentally different mechanism from the Fabry-P\'{e}rot cavity \citep{fabry1899}, which confines light spatially between mirrors. Here, confinement is topological: the network structure itself provides the return path.

\subsection{Scope}

We derive the following from the Bounded Phase Space Law (Axiom~\ref{ax:bpsl_recap}):
\begin{itemize}[nosep]
\item The absorption-emission interval as a physical tick (Theorem~\ref{thm:tick_emergence})
\item The nested tick hierarchy from emission lifetimes (Theorem~\ref{thm:hierarchical_nesting})
\item Harmonic proximity and network formation (Theorem~\ref{thm:harmonic_edge})
\item Closed loops and resonant circulation (Theorems~\ref{thm:loop_existence}--\ref{thm:virtual_cavity})
\item Self-clocking and self-validation (Theorems~\ref{thm:self_clocking}--\ref{thm:self_validation})
\item Entry from any node (Theorem~\ref{thm:entry_any_node})
\item Validation on six molecules against NIST data with zero adjustable parameters
\end{itemize}

No additional physical postulates are introduced beyond the Bounded Phase Space Law established in Paper~1.

\subsection{Paper Structure}

Part~I (Sections~\ref{sec:introduction}--\ref{sec:bps_recap}) establishes foundations, briefly recapitulating the bounded phase space framework. Part~II (Sections~\ref{sec:tick}--\ref{sec:hierarchy}) derives the molecular tick and nested tick hierarchy. Part~III (Sections~\ref{sec:harmonic_network}--\ref{sec:self_clocking}) constructs harmonic networks, proves loop existence, and establishes self-clocking and self-validation. Part~IV (Sections~\ref{sec:validation}--\ref{sec:error_analysis}) validates on six molecules. Part~V (Sections~\ref{sec:discussion}--\ref{sec:conclusion}) discusses implications and concludes.

%==============================================================================
\section{Bounded Phase Space Recap}
\label{sec:bps_recap}
%==============================================================================

We state the foundational results from Paper~1 \citep{sachikonye2025a} without repeating their proofs.

\subsection{The Axiom}

\begin{axiom}[The Bounded Phase Space Law]
\label{ax:bpsl_recap}
All persistent dynamical systems occupy bounded regions of phase space with finite Liouville measure, and these bounded regions admit hierarchical partitioning into distinguishable subregions.
\end{axiom}

A dynamical system is \emph{persistent} if it maintains its identity over time. The axiom requires that every such system is confined to a finite region $\Omega$ of phase space $\Gamma = \{(\mathbf{q}, \mathbf{p}) \in \Real^{2n}\}$ satisfying $\mu(\Omega) = \int_\Omega d^n q \, d^n p / h^n < \infty$, where $\mu$ denotes the Liouville measure \citep{planck1901, landau1976}.

\subsection{Oscillatory Necessity and Partition Coordinates}

\begin{theorem}[Oscillatory Necessity {\citep[Theorem~2.4]{sachikonye2025a}}]
\label{thm:oscillatory_recap}
For self-consistent dynamical systems in bounded phase space, oscillatory dynamics is the unique valid mode. Static, monotonic, and chaotic alternatives violate fundamental consistency requirements.
\end{theorem}

From the Poincar\'{e} recurrence theorem \citep{poincare1890} and bounded phase space geometry, the natural coordinate system is $(n, l, m, s)$ with shell capacity $C(n) = 2n^2$, and the frequency-energy identity is $E = \hbar\omega$ \citep{bohr1913, sachikonye2025a}.

\subsection{The Triple Equivalence Theorem}

\begin{theorem}[Triple Equivalence {\citep[Theorem~5.1]{sachikonye2025a}}]
\label{thm:triple_recap}
For any bounded oscillatory system, the oscillatory entropy, categorical entropy, and partition entropy are identical:
\begin{equation}
S_{\mathrm{osc}} = S_{\mathrm{cat}} = S_{\mathrm{part}} = \kB \Depth \ln n
\label{eq:triple_equivalence}
\end{equation}
where $\Depth$ is the partition depth and $n$ is the number of categorical states per partition level.
\end{theorem}

The Triple Equivalence establishes that oscillatory dynamics (counting oscillation cycles), categorical state enumeration (counting distinguishable states), and partition operations (performing hierarchical subdivision) are three descriptions of the \emph{same} mathematical object. A CPU clock oscillating at frequency $\omega$ performs $\omega/(2\pi)$ categorical state transitions per second. This identity, not analogy, is the foundation for what follows.

\subsection{The Commutation Theorem}

\begin{theorem}[Categorical-Physical Commutation {\citep[Theorem~8.1]{sachikonye2025a}}]
\label{thm:commutation_recap}
For any categorical observable $\hat{O}_{\mathrm{cat}}$ and any physical observable $\hat{O}_{\mathrm{phys}}$:
\begin{equation}
[\hat{O}_{\mathrm{cat}}, \hat{O}_{\mathrm{phys}}] = 0
\label{eq:commutation}
\end{equation}
Categorical measurement is quantum non-demolition: it determines partition coordinates without backaction.
\end{theorem}

\subsection{S-Entropy Coordinate Space}

The entropy coordinate space $\Sspace = [0,1]^3$ has coordinates $(\Sk, \St, \Se)$ representing knowledge, temporal, and evolution entropy respectively. Trajectories in $\Sspace$ represent the refinement of categorical distinctions over time. Trajectory completion---return to within $\epsilon$ of the starting point---constitutes one complete measurement cycle \citep{sachikonye2025a}.

%==============================================================================
% PART II: THE MOLECULAR TICK
%==============================================================================

\part{The Molecular Tick}

%==============================================================================
\section{The Absorption-Emission Interval as Physical Tick}
\label{sec:tick}
%==============================================================================

\subsection{Spontaneous Emission as Categorical Return}

Consider a molecule with ground state $|0\rangle$ and excited electronic state $|2\rangle$ separated by energy $\Delta E = \hbar\omega_0$. At time $t = 0$, the molecule absorbs a photon at frequency $\omega_0$ and transitions to $|2\rangle$. At time $t = \tem$, the molecule spontaneously emits a photon and returns to $|0\rangle$. The emission lifetime is \citep{einstein1917, dirac1928}:
\begin{equation}
\tem = \frac{1}{\Aki} = \frac{3\pi \epsilon_0 \hbar c^3}{\omega_0^3 |\langle 0 | \hat{\mu} | 2 \rangle|^2}
\label{eq:emission_lifetime}
\end{equation}
where $\Aki$ is the Einstein A coefficient for the transition $|2\rangle \to |0\rangle$, and $\langle 0 | \hat{\mu} | 2 \rangle$ is the transition dipole matrix element.

The interval $[0, \tem]$ represents one complete categorical cycle. The system begins in one partition cell ($|0\rangle$), transitions to another ($|2\rangle$), and returns to the first. This is precisely the Poincar\'{e} recurrence applied to the two-state partition $\{|0\rangle, |2\rangle\}$.

\begin{theorem}[Tick Emergence]
\label{thm:tick_emergence}
For a bounded oscillatory system with electronic transition at frequency $\omega_0$, the absorption-emission interval
\begin{equation}
\tem = \frac{1}{\Aki}
\label{eq:tick_period}
\end{equation}
where $\Aki$ is the Einstein A coefficient, constitutes a categorical tick: one complete cycle of the partition operation that distinguishes excited from ground state.
\end{theorem}

\begin{proof}
From the Bounded Phase Space Law (Axiom~\ref{ax:bpsl_recap}), the molecule occupies a bounded region $\Omega$ of phase space. The two electronic states $|0\rangle$ and $|2\rangle$ define a partition $\mathcal{P} = \{P_0, P_2\}$ of $\Omega$ into two subregions. By Oscillatory Necessity (Theorem~\ref{thm:oscillatory_recap}), the system undergoes oscillatory dynamics within $\Omega$.

The Poincar\'{e} recurrence theorem \citep{poincare1890} guarantees that the system returns to any subregion it visits. For the two-state partition $\mathcal{P}$, the mean return time from $P_0$ through $P_2$ and back to $P_0$ is the Poincar\'{e} return time for this partition. By Kac's lemma \citep{kac1947}, this return time is:
\begin{equation}
\langle \tau_{\mathrm{return}} \rangle = \frac{\mu(\Omega)}{\mu(P_0)}
\end{equation}

For a two-level system in thermal equilibrium at low temperature (where $P_0$ dominates), $\mu(P_0) \approx \mu(\Omega)$ and the return time is approximately one oscillation period. The physical return time for the specific transition $|0\rangle \to |2\rangle \to |0\rangle$ is precisely $\tem = 1/\Aki$, as determined by the Einstein A coefficient.

By the Triple Equivalence (Theorem~\ref{thm:triple_recap}), this oscillatory return period $\tem$ equals the categorical cycle time (one complete enumeration of the two-state partition) and equals the partition operation period (one complete application of the partition operation $\mathcal{P}$). Therefore $\tem$ constitutes a categorical tick: one complete cycle of the partition operation that distinguishes $|0\rangle$ from $|2\rangle$.
\end{proof}

\begin{remark}
The tick period $\tem$ is \emph{intrinsic} to the molecule. It is determined by three quantities: the transition frequency $\omega_0$, the transition dipole moment $|\langle 0 | \hat{\mu} | 2 \rangle|$, and the physical constants $\epsilon_0$, $\hbar$, $c$. All three are properties of the molecular partition structure. No external clock, no reference oscillator, no timing circuit is required. The molecule ticks itself.
\end{remark}

\subsection{The Tick Rate as Categorical Actualization Rate}

\begin{theorem}[Tick Rate Identity]
\label{thm:tick_rate}
The tick rate equals the categorical actualization rate:
\begin{equation}
\frac{d\Depth}{dt} = \frac{1}{\tem} = \Aki
\label{eq:tick_rate}
\end{equation}
where $\Depth$ is the partition depth (number of completed categorical cycles). This IS spectroscopic measurement: counting ticks IS measuring the molecule.
\end{theorem}

\begin{proof}
Each completed absorption-emission cycle increments the categorical state counter by one: $\Delta\Depth = 1$ per cycle. The time per cycle is $\tem$. Therefore:
\begin{equation}
\frac{d\Depth}{dt} = \frac{\Delta\Depth}{\tem} = \frac{1}{\tem} = \Aki
\end{equation}

By the Triple Equivalence, the categorical state count rate $d\Depth/dt$ is identical to the oscillation frequency $\omega_0/(2\pi)$ times the geometric factor encoding the partition structure. The Einstein A coefficient $\Aki$ encodes precisely this product: $\Aki \propto \omega_0^3 |\langle 0|\hat{\mu}|2\rangle|^2$. The factor $\omega_0^3$ arises from the density of photon states (the partition depth of the radiation field at frequency $\omega_0$), and $|\langle 0|\hat{\mu}|2\rangle|^2$ is the transition probability (the partition overlap between $|0\rangle$ and $|2\rangle$).

Therefore, measuring the Einstein A coefficient---which is standard spectroscopic practice---IS counting categorical ticks. The spectroscopist who measures an emission lifetime is directly measuring the categorical actualization rate.
\end{proof}

\subsection{The Absorption-Emission Cycle in S-Entropy Space}

The absorption-emission cycle traces a closed trajectory in the S-entropy coordinate space $\Sspace = [0,1]^3$:

\begin{definition}[Tick Trajectory]
\label{def:tick_trajectory}
For an absorption-emission cycle at frequency $\omega_0$ with emission lifetime $\tem$, the trajectory in $\Sspace$ is:
\begin{align}
\Sk(t) &= \begin{cases} \kB \ln(1 + \omega_0 t / \phi_0) & 0 \leq t < \tem/2 \\ \kB \ln(1 + \omega_0 (\tem - t)/\phi_0) & \tem/2 \leq t \leq \tem \end{cases} \label{eq:sk_trajectory} \\
\St(t) &= \kB \ln(1 + t/\tau_0) \label{eq:st_trajectory} \\
\Se(t) &= \begin{cases} \kB \ln(1 + \hbar\omega_0/E_0) & 0 < t < \tem \\ 0 & t = 0, \tem \end{cases} \label{eq:se_trajectory}
\end{align}
where $\phi_0$, $\tau_0$, $E_0$ are the normalization scales from Paper~1.
\end{definition}

The trajectory rises in $\Se$ upon absorption (energy increases), maintains the excited-state value during the excited-state lifetime, and returns to zero upon emission. The knowledge entropy $\Sk$ increases as the state is characterized and returns upon cycle completion. The temporal entropy $\St$ advances monotonically---time does not reverse. The key point is that $(\Sk, \Se)$ return to their initial values while $\St$ advances by the tick interval $\tem$. This is trajectory completion: the system's state-space coordinates return to their starting configuration, marking one categorical cycle.

\begin{figure*}[!htbp]
    \centering
    \includegraphics[width=\textwidth]{figures/panel_1_molecular_tick.png}
    \caption{\textbf{The Molecular Tick: Absorption-Emission as Categorical Cycle.}
    (\textbf{A})~Three-dimensional trajectory of a molecular absorption-emission cycle in S-entropy space $(\Sk, \St, \Se)$. The system begins at the ground state (bottom), spirals upward during photon absorption to the excited state (top torus), and decays back during spontaneous emission. The trajectory is coloured by time (yellow$\to$green$\to$blue), and the two circular orbits at ground and excited states represent the stationary categorical configurations. The emission lifetime $\tem$ marks one complete categorical cycle.
    (\textbf{B})~Emission lifetime versus vibrational frequency on logarithmic axes. The theoretical scaling $\tem \propto \omega^{-3}$ (dashed diagonal) follows from the Einstein A coefficient $A_{ki} = \omega_0^3 |\langle 0|\hat{\mu}|2\rangle|^2 / (3\pi\epsilon_0\hbar c^3)$. Points represent the fundamental modes of all six validation molecules, coloured by molecular type (diatomic$=$blue, triatomic$=$green, tetrahedral$=$orange, polyatomic$=$red). The spread about the theoretical line reflects variation in transition dipole moments.
    (\textbf{C})~Tick rate identity: categorical state transition rate $dM/dt = 1/\tem = A_{ki}$ versus frequency. The linear relationship on log-log axes confirms that the tick rate equals the Einstein A coefficient. Each point corresponds to one vibrational mode from the molecular database, with the diagonal reference line indicating exact proportionality.
    (\textbf{D})~Population dynamics during one emission cycle. Excited-state population $P_2(t) = e^{-t/\tem}$ (red, decaying) and ground-state population $P_0(t) = 1 - e^{-t/\tem}$ (blue, rising) versus normalised time $t/\tem$. The vertical dashed line at $t = \tem$ marks the emission lifetime---the point where the excited-state population has decayed to $1/e$. The crossing at $t = \tem \ln 2 \approx 0.69\,\tem$ marks equipartition between states. The interval $[0, \tem]$ constitutes one complete categorical tick.}
    \label{fig:molecular_tick}
\end{figure*}

%==============================================================================
\section{Nested Tick Hierarchy}
\label{sec:hierarchy}
%==============================================================================

\subsection{Multiple Emission Lifetimes}

A molecule possesses many oscillatory modes---electronic, vibrational, and rotational---each with its own emission lifetime. These lifetimes span orders of magnitude:

\begin{definition}[Emission Lifetime Hierarchy]
\label{def:lifetime_hierarchy}
For a molecule with $N$ oscillatory modes, the emission lifetimes form an ordered sequence:
\begin{equation}
\tem^{(1)} > \tem^{(2)} > \tem^{(3)} > \cdots > \tem^{(N)}
\label{eq:lifetime_ordering}
\end{equation}
Typical orders of magnitude are:
\begin{align}
\text{Rotational:} \quad & \tem^{(\mathrm{rot})} \sim 10^{-1} \text{--} 10^{3} \;\mathrm{s} \label{eq:rot_lifetime} \\
\text{Vibrational:} \quad & \tem^{(\mathrm{vib})} \sim 10^{-6} \text{--} 10^{-2} \;\mathrm{s} \label{eq:vib_lifetime} \\
\text{Electronic:} \quad & \tem^{(\mathrm{elec})} \sim 10^{-9} \text{--} 10^{-7} \;\mathrm{s} \label{eq:elec_lifetime}
\end{align}
\end{definition}

Within each category, different modes have different lifetimes. For a polyatomic molecule with $N_{\mathrm{vib}} = 3N_{\mathrm{atoms}} - 6$ vibrational modes (or $3N_{\mathrm{atoms}} - 5$ for linear molecules), each normal mode $\nu_i$ has a characteristic frequency $\omega_i$ and emission lifetime $\tem^{(i)}$. The lifetimes are determined by the transition dipole derivatives $\partial\mu/\partial Q_i$ where $Q_i$ is the normal coordinate \citep{wilson1955, herzberg1945}.

\subsection{Tree Structure from Tick Nesting}

\begin{theorem}[Hierarchical Nesting]
\label{thm:hierarchical_nesting}
For a molecule with $N$ oscillatory modes, the emission lifetimes $\{\tem^{(i)}\}_{i=1}^{N}$ form a hierarchically nested partition of the temporal axis. Each $\tem^{(i)}$ subdivides the interval of the next longer $\tem^{(j)}$ into $\lfloor \tem^{(j)}/\tem^{(i)} \rfloor$ complete categorical cycles.
\end{theorem}

\begin{proof}
Let $\tem^{(j)} > \tem^{(i)}$ for modes $j$ and $i$. During one complete cycle of mode $j$ (duration $\tem^{(j)}$), mode $i$ completes
\begin{equation}
N_{ij} = \left\lfloor \frac{\tem^{(j)}}{\tem^{(i)}} \right\rfloor
\label{eq:subdivision}
\end{equation}
categorical cycles. Each cycle of mode $i$ is a complete Poincar\'{e} return in the subspace spanned by mode $i$, and each of these returns is nested within the single cycle of mode $j$. By the Bounded Phase Space Law, the phase space region $\Omega_j$ associated with mode $j$ contains the subregion $\Omega_i$ associated with mode $i$ (since higher-frequency modes correspond to finer partition structure at greater depth). Therefore the categorical cycles of mode $i$ form a nested partition of the temporal interval $[0, \tem^{(j)}]$.

This nesting is hierarchical: if $\tem^{(k)} < \tem^{(i)} < \tem^{(j)}$, then mode $k$'s cycles subdivide mode $i$'s cycles, which subdivide mode $j$'s cycles. The result is a tree of nested temporal intervals.
\end{proof}

\begin{definition}[Tick Tree]
\label{def:tick_tree}
The \emph{tick tree} of a molecule is a rooted tree $T = (V, E_{\mathrm{tree}})$ where:
\begin{itemize}[nosep]
\item Vertices $V = \{v_1, v_2, \ldots, v_N\}$ correspond to the $N$ oscillatory modes
\item The root is the mode with the longest emission lifetime $\tem^{(1)}$
\item Vertex $v_i$ is a child of $v_j$ if $\tem^{(i)} < \tem^{(j)}$ and there is no intermediate mode $v_k$ with $\tem^{(i)} < \tem^{(k)} < \tem^{(j)}$ in the same branch
\item Each tree edge $(v_j, v_i) \in E_{\mathrm{tree}}$ is labeled with the subdivision number $N_{ij} = \lfloor \tem^{(j)}/\tem^{(i)} \rfloor$
\end{itemize}
\end{definition}

\begin{remark}
A single molecule may have multiple tick trees if it possesses independent sets of modes (e.g., one tree for the electronic-vibrational manifold and another for the rotational manifold). In practice, for polyatomic molecules, the electronic, vibrational, and rotational modes form a single hierarchical tree through the Born-Oppenheimer separation \citep{bornoppenheimer1927}: electronic transitions are the root, vibrational modes are children, and rotational modes are grandchildren.
\end{remark}

\subsection{Partition Depth from Tick Ratios}

\begin{theorem}[Partition Depth from Tick Ratio]
\label{thm:partition_depth_tick}
The partition depth between parent node $i$ and child node $j$ in the tick tree is:
\begin{equation}
\Depth_{ij} = \log_b\!\left(\frac{\tem^{(i)}}{\tem^{(j)}}\right)
\label{eq:partition_depth_tick}
\end{equation}
where $b$ is the partition base ($b = 3$ for ternary encoding in $\Sspace$). The total tree depth is:
\begin{equation}
\Depth_{\mathrm{total}} = \sum_{(i,j) \in E_{\mathrm{tree}}} \Depth_{ij} = \sum_{(i,j) \in E_{\mathrm{tree}}} \log_b\!\left(\frac{\tem^{(i)}}{\tem^{(j)}}\right)
\label{eq:total_tree_depth}
\end{equation}
\end{theorem}

\begin{proof}
The partition depth measures the number of refinement steps required to distinguish child from parent. Each refinement step subdivides the parent's temporal interval by a factor of $b$ (the partition base). To subdivide $\tem^{(i)}$ into intervals of length $\tem^{(j)}$, we require $\Depth_{ij}$ refinement steps where $b^{\Depth_{ij}} = \tem^{(i)}/\tem^{(j)}$, giving $\Depth_{ij} = \log_b(\tem^{(i)}/\tem^{(j)})$.

The total tree depth is the sum over all edges because each edge represents an independent sequence of refinement steps, and partition depths are additive along tree paths (each refinement is applied to the result of the previous one, so the total number of steps is the sum).
\end{proof}

\begin{corollary}[Entropy from Tick Hierarchy]
\label{cor:entropy_tick}
The categorical entropy stored in the tick tree is:
\begin{equation}
S_{\mathrm{tree}} = \kB \Depth_{\mathrm{total}} \ln b = \kB \sum_{(i,j) \in E_{\mathrm{tree}}} \ln\!\left(\frac{\tem^{(i)}}{\tem^{(j)}}\right)
\label{eq:tree_entropy}
\end{equation}
This equals the total information content of the molecular vibrational spectrum, measured in nats.
\end{corollary}

\begin{proof}
From the Triple Equivalence (Theorem~\ref{thm:triple_recap}), $S = \kB \Depth \ln n$. For base-$b$ partition, $n = b$ and $S = \kB \Depth \ln b$. Substituting $\Depth = \Depth_{\mathrm{total}}$ and using $\Depth_{ij} \ln b = \ln(\tem^{(i)}/\tem^{(j)})$ gives the result.
\end{proof}

\begin{figure*}[!htbp]
    \centering
    \includegraphics[width=\textwidth]{figures/panel_2_tick_hierarchy.png}
    \caption{\textbf{Nested Tick Hierarchy and Tree Structure.}
    (\textbf{A})~Three-dimensional tick tree for CH$_4$ (methane). Nodes represent vibrational modes plotted in (frequency, tree level, partition depth $\mathcal{M}$) space. The root node $\nu_4$ (1306~cm$^{-1}$) at level 0 connects downward to $\nu_2$ (1534~cm$^{-1}$) at level 1, which branches to $\nu_1$ (2917~cm$^{-1}$) and $\nu_3$ (3019~cm$^{-1}$) at level 2. Node size scales with subdivision number $N_{ij} = \lfloor\omega_j/\omega_i\rfloor$, and colour encodes frequency (blue$=$low, yellow$=$high). Solid lines are tree edges connecting parent to child.
    (\textbf{B})~Subdivision numbers $N_{ij}$ for all parent$\to$child tree edges across all six molecules. Each bar represents one edge, grouped and coloured by molecule. Most subdivision numbers are 1--3, reflecting the typical 2--3$\times$ frequency ratios between adjacent vibrational modes. The tall bar for H$_2$ ($N = 37$) reflects the large ratio between its vibrational and rotational frequencies.
    (\textbf{C})~Partition depth $\mathcal{M}_{ij} = \log_3(\omega_j/\omega_i)$ versus frequency ratio $\omega_j/\omega_i$ for all tree edges. The smooth $\log_3$ reference curve (solid line) passes through all data points, confirming that partition depth is the base-3 logarithm of the frequency ratio. Points are coloured by molecule. Higher frequency ratios produce deeper partition nesting.
    (\textbf{D})~Frequency ordering of all vibrational modes across the six molecules, displayed as horizontal bars sorted by frequency within each molecular group. Bar length is proportional to mode frequency in cm$^{-1}$, and colour identifies the molecule. This ordering defines the tree structure: within each molecule, the lowest-frequency mode is the tree root and higher-frequency modes are descendants. The frequency span from CCl$_4$ ($\nu_2 = 218$~cm$^{-1}$) to H$_2$ (4401~cm$^{-1}$) covers a factor of 20.}
    \label{fig:tick_hierarchy}
\end{figure*}

\subsection{Example: The Tick Tree of Water}

Water (H$_2$O) has three vibrational modes: $\nu_1 = 3657$ cm$^{-1}$ (symmetric stretch), $\nu_2 = 1595$ cm$^{-1}$ (bend), and $\nu_3 = 3756$ cm$^{-1}$ (asymmetric stretch) \citep{shimanouchi1972, nist_webbook}. Converting to angular frequencies $\omega_i = 2\pi c \nu_i$:
\begin{align}
\omega_1 &= 2\pi \times (3.00 \times 10^{10}) \times 3657 = 6.883 \times 10^{14} \;\mathrm{rad/s} \\
\omega_2 &= 2\pi \times (3.00 \times 10^{10}) \times 1595 = 3.003 \times 10^{14} \;\mathrm{rad/s} \\
\omega_3 &= 2\pi \times (3.00 \times 10^{10}) \times 3756 = 7.070 \times 10^{14} \;\mathrm{rad/s}
\end{align}

The vibrational emission lifetimes scale as $\tem \propto \omega^{-3} |\partial\mu/\partial Q|^{-2}$. For water, the asymmetric stretch $\nu_3$ has the largest transition dipole derivative and the highest frequency, giving the shortest $\tem$. The bend $\nu_2$ has the longest $\tem$ (lowest frequency, moderate dipole derivative). The ordering is:
\begin{equation}
\tem^{(\nu_2)} > \tem^{(\nu_1)} > \tem^{(\nu_3)}
\end{equation}

The tick tree has $\nu_2$ as root, with $\nu_1$ and $\nu_3$ as children. The subdivision numbers are:
\begin{equation}
N_{21} = \left\lfloor \frac{\omega_1}{\omega_2} \right\rfloor = \left\lfloor \frac{3657}{1595} \right\rfloor = 2, \qquad
N_{23} = \left\lfloor \frac{\omega_3}{\omega_2} \right\rfloor = \left\lfloor \frac{3756}{1595} \right\rfloor = 2
\end{equation}

Each bend cycle contains approximately 2 complete symmetric stretch cycles and 2 complete asymmetric stretch cycles. The partition depth from bend to symmetric stretch is $\Depth_{21} = \log_3(3657/1595) = \log_3(2.293) = 0.756$.

%==============================================================================
% PART III: FROM TREES TO NETWORKS
%==============================================================================

\part{From Trees to Networks}

%==============================================================================
\section{Harmonic Proximity and Network Formation}
\label{sec:harmonic_network}
%==============================================================================

\subsection{Harmonic Proximity}

Trees are hierarchical: each node has exactly one parent (except the root). The topology is acyclic---there are no closed loops. To obtain loops, which are necessary for light circulation, we must add edges that connect nodes across different branches or different trees. The criterion for adding such edges is harmonic proximity.

\begin{definition}[Harmonic Proximity]
\label{def:harmonic_proximity}
Two oscillatory modes with frequencies $\omega_i$ and $\omega_j$ are \emph{harmonically proximate} if their frequency ratio is a low-order rational number:
\begin{equation}
\frac{\omega_i}{\omega_j} = \frac{p}{q}, \qquad p, q \in \Natural^+, \quad \gcd(p, q) = 1, \quad \max(p, q) \leq \qmax
\label{eq:harmonic_proximity}
\end{equation}
The \emph{harmonic proximity order} is $\eta = \max(p, q)$. Lower $\eta$ means stronger harmonic coupling. The threshold $\qmax$ (typically $\qmax = 10$) defines the maximum harmonic order considered.
\end{definition}

\begin{remark}
Harmonic proximity is the number-theoretic structure underlying resonance. Two modes with $\omega_i/\omega_j = 2/1$ are in perfect octave resonance; $3/2$ is a perfect fifth; $4/3$ is a perfect fourth. The connection to musical harmony is not accidental---both molecular and acoustic resonance are consequences of bounded oscillatory dynamics \citep{strogatz2015}. The rational approximations to frequency ratios are systematically enumerated by Farey sequences and the Stern-Brocot tree \citep{hardy1979}.
\end{remark}

In practice, exact rational ratios are never achieved---the ratio $\omega_i/\omega_j$ is an irrational number that can be approximated by $p/q$ to some precision. We define the \emph{harmonic deviation}:

\begin{definition}[Harmonic Deviation]
\label{def:harmonic_deviation}
The harmonic deviation of the pair $(\omega_i, \omega_j)$ from the rational approximation $p/q$ is:
\begin{equation}
\delta_{ij} = \left| \frac{\omega_i}{\omega_j} - \frac{p}{q} \right|
\label{eq:harmonic_deviation}
\end{equation}
where $p/q$ is the best rational approximation with $\max(p, q) \leq \qmax$. Two modes are considered harmonically proximate if $\delta_{ij} < \delta_{\max}$, where $\delta_{\max}$ is the anharmonicity threshold (we use $\delta_{\max} = 0.05$ throughout).
\end{definition}

\subsection{Harmonic Edge Formation}

\begin{theorem}[Harmonic Edge Formation]
\label{thm:harmonic_edge}
For two modes $i, j$ on different branches of the tick tree (or on different trees), if they are harmonically proximate with order $\eta \leq \etamax$, they share a harmonic edge with coupling strength:
\begin{equation}
g_{ij} = \frac{1}{\eta^2} \cdot \frac{|\langle \psi_i | \hat{\mu} | \psi_j \rangle|^2}{\hbar^2 \omega_i \omega_j}
\label{eq:coupling_strength}
\end{equation}
where $|\langle \psi_i | \hat{\mu} | \psi_j \rangle|$ is the transition dipole matrix element between modes $i$ and $j$.
\end{theorem}

\begin{proof}
Consider two oscillatory modes $i$ and $j$ with frequencies $\omega_i$ and $\omega_j$ satisfying $\omega_i/\omega_j \approx p/q$. By the frequency-energy identity $E = \hbar\omega$, the energies satisfy $E_i/E_j \approx p/q$. A rational energy ratio means that $q$ quanta of mode $i$ have the same energy as $p$ quanta of mode $j$:
\begin{equation}
q \cdot \hbar\omega_i \approx p \cdot \hbar\omega_j
\label{eq:resonance_condition}
\end{equation}

This is the condition for Fermi resonance \citep{fermi1931}: energy exchange between modes $i$ and $j$ is resonantly enhanced when their energies are commensurately related. The coupling strength for Fermi resonance is proportional to the transition dipole matrix element and inversely proportional to the order of the resonance \citep{herzberg1945, wilson1955}.

The factor $1/\eta^2$ encodes the decrease in coupling strength with resonance order: higher-order resonances ($\eta = \max(p,q)$ large) require more quanta to be exchanged simultaneously, and the probability of multi-quantum exchange decreases as $\eta^{-2}$ (from the perturbation theory matrix element at order $\eta$) \citep{landau1977}. The factor $|\langle \psi_i|\hat{\mu}|\psi_j\rangle|^2/(\hbar^2\omega_i\omega_j)$ is the dimensionless transition probability between modes, normalized by their zero-point energies.
\end{proof}

\begin{remark}
The coupling $g_{ij}$ is not an adjustable parameter---it is determined entirely by the molecular partition structure (transition dipole moments, frequencies). The only threshold parameter is $\etamax$, which sets the maximum harmonic order considered. Setting $\etamax = 10$ captures all physically significant resonances while excluding negligibly weak high-order couplings.
\end{remark}

\subsection{The Harmonic Molecular Network}

\begin{definition}[Harmonic Molecular Network]
\label{def:harmonic_network}
The \emph{harmonic molecular network} is a graph $G = (V, E)$ where:
\begin{itemize}[nosep]
\item Vertices $V = \{v_1, v_2, \ldots, v_N\}$ are all oscillatory modes of the molecule (electronic, vibrational, rotational)
\item Tree edges $E_{\mathrm{tree}}$: parent-child relationships from the tick hierarchy (Section~\ref{sec:hierarchy})
\item Harmonic edges $E_{\mathrm{harm}}$: harmonic proximity connections across different branches or trees
\item $E = E_{\mathrm{tree}} \cup E_{\mathrm{harm}}$
\end{itemize}
Each tree edge $(v_j, v_i) \in E_{\mathrm{tree}}$ is labeled with the subdivision number $N_{ij}$. Each harmonic edge $(v_i, v_j) \in E_{\mathrm{harm}}$ is labeled with the rational approximation $p/q$, the harmonic order $\eta$, and the coupling strength $g_{ij}$.
\end{definition}

\begin{theorem}[Network Connectivity]
\label{thm:network_connectivity}
For any molecule with $N_{\mathrm{vib}} \geq 3$ vibrational modes, the harmonic molecular network is connected.
\end{theorem}

\begin{proof}
We prove connectivity in two steps.

\textbf{Step 1: Tree connectivity.} The tick tree is connected by definition (it is a tree). Therefore all modes within a single tick tree are connected through tree edges.

\textbf{Step 2: Harmonic edges ensure cross-branch connectivity.} For a polyatomic molecule with $N_{\mathrm{vib}} \geq 3$ vibrational modes, the vibrational frequencies $\{\omega_i\}$ are determined by the eigenvalues of the molecular force constant matrix $\mathbf{F}$ \citep{wilson1955}. The force constant matrix couples all internal coordinates through shared atoms: if atoms A and B participate in mode $i$, and atoms B and C participate in mode $j$, then modes $i$ and $j$ share the force constant elements involving atom B. This coupling ensures that the frequency ratios $\omega_i/\omega_j$ span a continuous range.

By the three-distance theorem \citep{hardy1979}, for any three frequencies $\omega_1, \omega_2, \omega_3$ that are not all pairwise commensurable, at least one pair $(\omega_i, \omega_j)$ satisfies $|\omega_i/\omega_j - p/q| < 1/(2q^2)$ for some $p/q$ with $q \leq \qmax$. Since the vibrational frequencies of a polyatomic molecule are generically incommensurable (determined by distinct eigenvalues of the force constant matrix), at least one harmonic edge exists between any two branches of the tick tree.

Therefore the combined graph $G = (V, E_{\mathrm{tree}} \cup E_{\mathrm{harm}})$ is connected.
\end{proof}

\begin{figure*}[!htbp]
    \centering
    \includegraphics[width=\textwidth]{figures/panel_3_harmonic_network.png}
    \caption{\textbf{Harmonic Network Formation.}
    (\textbf{A})~Three-dimensional harmonic molecular network for C$_6$H$_6$ (benzene). Five vibrational modes are plotted as nodes in (frequency, mode index, coupling) space, with node size proportional to frequency. Solid blue lines represent tree edges (parent-child hierarchy); dashed red lines represent harmonic edges (cross-branch connections where $\omega_i/\omega_j \approx p/q$). The harmonic edges transform the tree into a network admitting closed loops. Notable ratios include $\nu_3/\nu_4 = 1178/673 = 1.750 \approx 7/4$ ($\delta = 0.001$) and $\nu_5/\nu_4 = 993/673 = 1.475 \approx 3/2$ ($\delta = 0.025$).
    (\textbf{B})~Distribution of harmonic ratio deviations $\delta = |\omega_i/\omega_j - p/q|$ for all 22 harmonic edges across the six molecules. Orange bars indicate tree-overlap edges (harmonic edge coincident with an existing tree edge); blue bars indicate cross-branch edges (novel connections). The vertical red line marks the tolerance threshold $\delta_{\max} = 0.05$. All edges satisfy $\delta < \delta_{\max}$, with the median deviation at $\delta \approx 0.02$.
    (\textbf{C})~Rational approximation quality: deviation $\delta$ versus observed frequency ratio $\omega_i/\omega_j$ for all harmonic edges. Point size is inversely proportional to the harmonic order $\eta = \max(p,q)$, so lower-order (stronger) harmonics appear as larger points. Colour encodes molecule identity. The horizontal red line at $\delta = 0.05$ separates accepted edges (below) from rejected pairs (above). The tightest harmonic coincidence ($\delta = 0.001$) belongs to benzene's $\nu_3/\nu_4 \approx 7/4$.
    (\textbf{D})~Network complexity scaling: number of harmonic edges $|E_{\mathrm{harm}}|$ versus number of vibrational modes $N$. Points are coloured and labelled by molecule. The grey curve shows the combinatorial upper bound $N(N-1)/2$ (all pairs harmonic). Actual edge counts fall well below this bound, indicating that harmonic proximity is a selective property. The growth from 1 edge (H$_2$, 2 modes) to 6 edges (CH$_4$ and C$_6$H$_6$, 4--5 modes) demonstrates the quadratic scaling potential.}
    \label{fig:harmonic_network}
\end{figure*}

%==============================================================================
\section{Closed Loops and Light Circulation}
\label{sec:loops}
%==============================================================================

\subsection{Loop Existence}

This section contains the central results of the paper.

\begin{theorem}[Loop Existence]
\label{thm:loop_existence}
For any harmonic molecular network with $|E_{\mathrm{harm}}| > 0$, there exists at least one closed loop traversing both tree edges and harmonic edges.
\end{theorem}

\begin{proof}
A tree on $N$ vertices has exactly $N - 1$ edges and is acyclic: every pair of vertices is connected by a unique path. Adding any edge to a tree creates exactly one cycle---this is a fundamental theorem of graph theory \citep{diestel2017}.

Let $(v_i, v_j) \in E_{\mathrm{harm}}$ be a harmonic edge. Since the tick tree is connected, there exists a unique tree path $P_{ij}$ from $v_i$ to $v_j$ traversing tree edges. The harmonic edge $(v_i, v_j)$ together with the tree path $P_{ij}$ forms a closed loop:
\begin{equation}
L = (v_i, \underbrace{v_{i_1}, v_{i_2}, \ldots, v_{i_k}}_{\text{tree path } P_{ij}}, v_j, v_i)
\label{eq:loop}
\end{equation}
where the last edge $(v_j, v_i)$ is the harmonic edge. This loop contains at least one tree edge (from $P_{ij}$) and exactly one harmonic edge, satisfying the requirements.

Since $|E_{\mathrm{harm}}| > 0$ by hypothesis, at least one such loop exists.
\end{proof}

\begin{corollary}[Cycle Rank]
\label{cor:cycle_rank}
The number of independent loops in the harmonic molecular network is:
\begin{equation}
C = |E| - |V| + 1 = |E_{\mathrm{tree}}| + |E_{\mathrm{harm}}| - |V| + 1 = |E_{\mathrm{harm}}|
\label{eq:cycle_rank}
\end{equation}
since $|E_{\mathrm{tree}}| = |V| - 1$ for a tree. Each harmonic edge creates exactly one independent loop.
\end{corollary}

\subsection{Circulation}

\begin{definition}[Circulation]
\label{def:circulation}
A \emph{circulation} is a closed loop $L = (v_1, v_2, \ldots, v_k, v_1)$ in the harmonic molecular network such that:
\begin{enumerate}[nosep]
\item The loop contains at least one tree edge and at least one harmonic edge.
\item Each edge traversal corresponds to a physical process:
\begin{itemize}[nosep]
\item Tree edge (parent $\to$ child): spontaneous emission at frequency $\omega_{\mathrm{parent}} - \omega_{\mathrm{child}}$
\item Tree edge (child $\to$ parent): absorption at frequency $\omega_{\mathrm{parent}} - \omega_{\mathrm{child}}$
\item Harmonic edge: energy exchange through Fermi resonance at the beat frequency
\end{itemize}
\item The total phase accumulated around the loop satisfies the constructive interference condition:
\begin{equation}
\Phi_L = \sum_{(i,j) \in L} \phi_{ij} = 2\pi n, \qquad n \in \Integer^+
\label{eq:constructive_interference}
\end{equation}
\end{enumerate}
\end{definition}

\begin{theorem}[Resonant Circulation]
\label{thm:resonant_circulation}
The circulation period of loop $L$ is:
\begin{equation}
T_L = \sum_{(i,j) \in L} \tau_{ij}
\label{eq:circulation_period}
\end{equation}
where $\tau_{ij}$ is the transit time for edge $(i,j)$:
\begin{align}
\tau_{ij} &= \tem^{(i)} && \text{for tree edges (emission/absorption)} \label{eq:tree_transit} \\
\tau_{ij} &= \frac{2\pi}{|\omega_i - \omega_j|} = \frac{1}{|\nu_i - \nu_j| c} && \text{for harmonic edges (beat frequency)} \label{eq:harmonic_transit}
\end{align}
The phase accumulated along each edge is:
\begin{equation}
\phi_{ij} = \omega_i \cdot \tau_{ij}
\label{eq:phase_per_edge}
\end{equation}
\end{theorem}

\begin{proof}
\textbf{Tree edge transit time.} A tree edge from parent $v_i$ to child $v_j$ represents a transition between the two modes. The transit time is the emission lifetime $\tem^{(i)}$ of the parent mode, since the parent must complete one absorption-emission cycle to transfer energy to the child. By Theorem~\ref{thm:tick_emergence}, this is one categorical tick of the parent.

\textbf{Harmonic edge transit time.} A harmonic edge between modes $v_i$ and $v_j$ represents energy exchange through Fermi resonance. The exchange occurs at the beat frequency $|\omega_i - \omega_j|$. The time for one complete beat cycle (full energy transfer and return) is $\tau_{ij} = 2\pi/|\omega_i - \omega_j|$.

\textbf{Phase accumulation.} During transit through edge $(i,j)$, mode $i$ oscillates at frequency $\omega_i$ for time $\tau_{ij}$, accumulating phase $\phi_{ij} = \omega_i \tau_{ij}$. The total phase around the loop is $\Phi_L = \sum \phi_{ij}$.

\textbf{Circulation period.} The total time for one traversal of the loop is the sum of all transit times. This follows from the sequential nature of the physical processes: emission must complete before the next absorption or resonance transfer can begin.
\end{proof}

\subsection{Virtual Resonant Cavities}

\begin{theorem}[Virtual Resonant Cavity]
\label{thm:virtual_cavity}
A closed loop $L$ satisfying the constructive interference condition (Equation~\ref{eq:constructive_interference}) constitutes a virtual resonant cavity with the following properties:
\begin{enumerate}[nosep]
\item Light circulates within the loop without walls, mirrors, or spatial confinement.
\item Confinement is categorical: the loop closes because partition coordinates return to initial values.
\item The cavity quality factor is:
\begin{equation}
Q_L = \frac{\bar{\omega} \, T_L}{2\pi}
\label{eq:quality_factor}
\end{equation}
where $\bar{\omega} = (1/|L|) \sum_{v_i \in L} \omega_i$ is the mean frequency of the modes in the loop.
\item There is no photon propagation bottleneck: coupling is categorical, not spatial.
\end{enumerate}
\end{theorem}

\begin{proof}
\textbf{(1) Circulation without walls.} In a conventional Fabry-P\'{e}rot cavity \citep{fabry1899, siegman1986}, light bounces between mirrors separated by distance $d$. The mirrors provide spatial confinement: the photon is reflected back. In the harmonic molecular network, the ``reflection'' is provided by the network topology: when light reaches a node and the only available continuation is through a harmonic edge back toward the starting node, the topology itself redirects the energy. No spatial mirrors are needed.

Formally, the circulation is a sequence of energy transfers:
\begin{equation}
v_1 \xrightarrow{\text{emit}} v_2 \xrightarrow{\text{absorb}} v_3 \xrightarrow{\text{resonate}} v_4 \xrightarrow{\text{emit}} \cdots \xrightarrow{} v_1
\end{equation}
Each arrow represents a physical process (emission, absorption, or Fermi resonance). The loop closes because the sequence of processes returns the system to its initial state.

\textbf{(2) Categorical confinement.} At each node $v_i$, the system occupies a definite partition state $(n_i, l_i, m_i, s_i)$. Traversing the loop transforms the partition state through a sequence of operations. The constructive interference condition $\Phi_L = 2\pi n$ ensures that after one complete loop, the partition state returns to its initial value. This is a topological constraint, not a spatial one.

\textbf{(3) Quality factor.} The quality factor of a resonant cavity is $Q = \omega_0 \tau_{\mathrm{lifetime}}$, where $\tau_{\mathrm{lifetime}}$ is the time for the field energy to decay to $1/e$ of its initial value \citep{siegman1986, vahala2003}. For the virtual cavity, the ``field lifetime'' is the circulation period $T_L$ (after one complete circulation, the system returns to its initial state and begins a new cycle). The mean frequency $\bar{\omega}$ serves as $\omega_0$.

\textbf{(4) No propagation bottleneck.} In a spatial cavity, the round-trip time is limited by the speed of light: $T_{\mathrm{RT}} = 2d/c$. The virtual cavity has no such limitation. Energy transfer between harmonically coupled modes occurs through categorical coupling---the transition matrix element $\langle \psi_i | \hat{\mu} | \psi_j \rangle$---which does not require photon propagation through space. The coupling is ``instantaneous'' in the sense that it is limited by the beat frequency $|\omega_i - \omega_j|$, not by $d/c$.
\end{proof}

\begin{remark}
The virtual resonant cavity differs from conventional cavities in a fundamental way: it has no spatial boundaries. A Fabry-P\'{e}rot cavity requires mirrors; a whispering gallery mode requires a curved surface \citep{vahala2003}; a photonic crystal cavity requires a periodic dielectric structure \citep{noda2007}. The virtual cavity requires only the network topology of harmonically coupled oscillatory modes. The cavity exists because the mathematics closes, not because the geometry does.
\end{remark}

\subsection{Constructive Interference Condition}

\begin{proposition}[Constructive Interference for Harmonic Loops]
\label{prop:constructive}
For a loop $L$ containing modes with frequencies $\omega_i / \omega_j = p/q$ (exact rational ratios), the constructive interference condition $\Phi_L = 2\pi n$ is automatically satisfied.
\end{proposition}

\begin{proof}
Consider a loop with two modes $\omega_i$ and $\omega_j$ connected by a harmonic edge with $\omega_i/\omega_j = p/q$ and a tree path. The total phase is:
\begin{align}
\Phi_L &= \omega_i \cdot \tem^{(i)} + \omega_j \cdot \tem^{(j)} + \omega_i \cdot \frac{2\pi}{|\omega_i - \omega_j|}
\end{align}

For the tree edges: $\omega_i \cdot \tem^{(i)} = \omega_i / \Aki^{(i)}$, which is a property of the transition. The emission lifetime is related to the frequency by $\Aki \propto \omega^3$, so $\omega \cdot \tem = \omega / \Aki \propto \omega^{-2}$.

For the harmonic edge: if $\omega_i/\omega_j = p/q$ exactly, then $|\omega_i - \omega_j| = \omega_j |p/q - 1| = \omega_j (p-q)/q$ (assuming $p > q$). The phase accumulated is:
\begin{equation}
\omega_i \cdot \frac{2\pi}{|\omega_i - \omega_j|} = \frac{p}{q} \cdot \omega_j \cdot \frac{2\pi q}{\omega_j(p - q)} = \frac{2\pi p}{p - q}
\end{equation}
which is $2\pi$ times a rational number. Summing over all edges in the loop, the total phase is $2\pi$ times a rational number, and for loops satisfying the harmonic condition at all edges, this rational number is an integer.
\end{proof}

\begin{remark}
For realistic molecules, the frequency ratios are not exactly rational but are rational to within the anharmonicity tolerance $\delta_{\max}$. The deviation from exact constructive interference is:
\begin{equation}
|\Phi_L - 2\pi n| \leq 2\pi \sum_{(i,j) \in L} \delta_{ij} \cdot \frac{\omega_j}{\omega_j - \omega_i} \cdot \frac{\omega_i}{\Aki^{(i)}}
\end{equation}
For $\delta_{ij} < 0.05$ and typical molecular parameters, this deviation is less than $0.1 \times 2\pi$, well within the linewidth of the constituent transitions.
\end{remark}

\begin{figure*}[!htbp]
    \centering
    \includegraphics[width=\textwidth]{figures/panel_4_circulation.png}
    \caption{\textbf{Closed Loops and Light Circulation.}
    (\textbf{A})~Three-dimensional loop visualisation for H$_2$O. The three vibrational modes ($\nu_1 = 3657$, $\nu_2 = 1595$, $\nu_3 = 3756$~cm$^{-1}$) form nodes in frequency space, connected by two tree edges (blue solid: $\nu_1 \to \nu_2$, $\nu_2 \to \nu_3$) and one harmonic edge (red dashed: $\nu_3 \to \nu_1$, $\omega_3/\omega_1 = 1.027 \approx 1/1$). The closed loop $\nu_1 \to \nu_2 \to \nu_3 \to \nu_1$ constitutes a virtual resonant cavity: light circulates through the loop via categorical coupling, not spatial confinement. The $z$-axis encodes the beat frequency $\Delta\omega$ for each edge.
    (\textbf{B})~Circulation periods $T_L$ for all 11 cross-branch loops across the six molecules, displayed on a logarithmic scale. Bars are coloured by molecule. Periods span three orders of magnitude: from $\sim$0.02~ps (CH$_4$, tight loops between closely spaced modes) to $\sim$9~ps (CO rovibrational loops with narrow beat frequencies). Each period is determined entirely by the molecular partition structure with zero adjustable parameters.
    (\textbf{C})~Beat frequency spectrum: vertical lines at each harmonic edge beat frequency $\Delta\omega = |\omega_i - \omega_j|$ in cm$^{-1}$, coloured by molecule. The spectrum spans from 99~cm$^{-1}$ (H$_2$O, $\nu_3 - \nu_1$) to 1713~cm$^{-1}$ (CH$_4$, $\nu_3 - \nu_4$). The CO$_2$ Fermi resonance at $\Delta\omega = 961$~cm$^{-1}$ ($\nu_3 - \nu_1$) is visible as a characteristic signature of the $2\nu_2 \approx \nu_1$ near-degeneracy first identified by \citet{fermi1931}.
    (\textbf{D})~Cumulative phase accumulation $\phi/(2\pi)$ around closed loops for H$_2$O (green) and CO$_2$ (red). The horizontal axis tracks normalised loop position from 0 (start) to 1 (return). Constructive interference requires the total phase to equal $2\pi n$ for integer $n$ (horizontal gridlines). Both curves reach integer multiples of $2\pi$ at loop completion, confirming the resonant circulation condition. The slope at each segment is proportional to the local frequency.}
    \label{fig:circulation}
\end{figure*}

%==============================================================================
\section{Self-Clocking and Self-Validation}
\label{sec:self_clocking}
%==============================================================================

\subsection{Self-Clocking}

\begin{theorem}[Self-Clocking]
\label{thm:self_clocking}
Each complete circulation around a resonant loop constitutes one tick of a self-contained clock. The circulation IS the tick; no external reference is needed.
\end{theorem}

\begin{proof}
The circulation period $T_L$ is determined entirely by the molecular partition structure:
\begin{equation}
T_L = \sum_{(i,j) \in L} \tau_{ij}
\end{equation}
where each $\tau_{ij}$ depends only on molecular frequencies and Einstein A coefficients---intrinsic molecular properties.

Each completion of the loop returns the system to its initial categorical state. By definition, a closed loop in the harmonic network begins and ends at the same vertex with the same partition coordinates. Therefore, the number of completed circulations at time $t$ is:
\begin{equation}
N_{\mathrm{circ}}(t) = \left\lfloor \frac{t}{T_L} \right\rfloor
\end{equation}
and the categorical actualization rate is:
\begin{equation}
\frac{d\Depth}{dt} = \frac{1}{T_L}
\end{equation}

By the Triple Equivalence (Theorem~\ref{thm:triple_recap}), this categorical state counting IS oscillation IS partition operation. The circulation provides all three simultaneously:
\begin{enumerate}[nosep]
\item \textbf{Oscillation:} The periodic return to the initial state is an oscillation with period $T_L$.
\item \textbf{State counting:} Each completed loop increments the categorical state counter.
\item \textbf{Partition operation:} Each edge traversal performs a partition operation (refining or coarsening the categorical description).
\end{enumerate}

No external clock is required because the circulation defines time operationally: one tick IS one circulation. The molecule is its own clock.
\end{proof}

\begin{corollary}[Circulation Frequency]
\label{cor:circulation_frequency}
The circulation defines a characteristic frequency:
\begin{equation}
\omega_L = \frac{2\pi}{T_L}
\end{equation}
and a characteristic energy:
\begin{equation}
E_L = \hbar\omega_L = \frac{2\pi\hbar}{T_L} = \frac{h}{T_L}
\end{equation}
This energy is the ``cost'' of one complete categorical cycle around the loop.
\end{corollary}

\subsection{Self-Validation}

\begin{theorem}[Self-Validation through Loop Redundancy]
\label{thm:self_validation}
If the harmonic molecular network contains $C > 1$ independent loops sharing common edges, then measurement of any partition coordinate through one loop can be cross-validated through every other loop traversing the same node.
\end{theorem}

\begin{proof}
Let $L_\alpha$ and $L_\beta$ be two independent loops sharing a common vertex $v_k$. Through loop $L_\alpha$, the partition coordinates of $v_k$ are determined by:
\begin{equation}
(n_k, l_k, m_k, s_k)_\alpha = f_\alpha(\{\omega_i, \Aki^{(i)}\}_{v_i \in L_\alpha})
\end{equation}
where $f_\alpha$ is the function that extracts partition coordinates from the loop parameters. Through loop $L_\beta$:
\begin{equation}
(n_k, l_k, m_k, s_k)_\beta = f_\beta(\{\omega_j, \Aki^{(j)}\}_{v_j \in L_\beta})
\end{equation}

Since $L_\alpha$ and $L_\beta$ are independent loops (different edges), the functions $f_\alpha$ and $f_\beta$ use different input data. Yet both must yield the same partition coordinates for the shared vertex $v_k$, because $v_k$'s partition coordinates are intrinsic to the molecular mode---they do not depend on the measurement path.

Formally, this follows from the Commutation Theorem (Theorem~\ref{thm:commutation_recap}): $[\hat{O}_{\mathrm{cat}}, \hat{O}_{\mathrm{phys}}] = 0$ guarantees that the partition coordinate assignment is independent of the physical path used to determine it. The categorical observable $\hat{O}_{\mathrm{cat}}$ (which assigns partition coordinates to $v_k$) commutes with all physical observables, including the choice of measurement loop. Therefore:
\begin{equation}
(n_k, l_k, m_k, s_k)_\alpha = (n_k, l_k, m_k, s_k)_\beta
\end{equation}
for all pairs of loops sharing vertex $v_k$.
\end{proof}

\begin{corollary}[Validation Multiplicity]
\label{cor:validation_multiplicity}
For a molecule with $C$ independent loops (Corollary~\ref{cor:cycle_rank}), the self-validation multiplicity is $C$. Each additional loop beyond the first provides an independent measurement channel for every vertex it traverses. The total number of cross-checks is:
\begin{equation}
N_{\mathrm{check}} = \sum_{v \in V} \left[ \binom{d_{\mathrm{loop}}(v)}{2} \right]
\end{equation}
where $d_{\mathrm{loop}}(v)$ is the number of independent loops passing through vertex $v$.
\end{corollary}

\subsection{Entry from Any Node}

\begin{theorem}[Entry from Any Node]
\label{thm:entry_any_node}
The network topology permits entry from any node. Any oscillatory mode of the molecule can serve as the starting point for circulation, and the measured partition coordinates are independent of the entry point.
\end{theorem}

\begin{proof}
\textbf{Step 1: Reachability.} The harmonic molecular network is connected (Theorem~\ref{thm:network_connectivity}). Therefore, for any vertex $v_i$ and any loop $L$, there exists a path from $v_i$ to some vertex $v_k \in L$. Once $v_k$ is reached, the circulation around $L$ can proceed.

\textbf{Step 2: Entry-point independence.} Consider a loop $L = (v_1, v_2, \ldots, v_k, v_1)$. The same loop can be traversed starting from any vertex:
\begin{align}
L_1 &= (v_1, v_2, \ldots, v_k, v_1) \\
L_2 &= (v_2, v_3, \ldots, v_k, v_1, v_2) \\
&\;\;\vdots \\
L_k &= (v_k, v_1, v_2, \ldots, v_{k-1}, v_k)
\end{align}
These are cyclic permutations of the same loop. The total phase $\Phi_L$, the circulation period $T_L$, and the partition coordinates assigned to each node are invariant under cyclic permutation.

\textbf{Step 3: Independence from entry path.} The path from the entry node $v_i$ to the loop $L$ is a tree path (since the tree connects all nodes). The tree path traversal performs partition operations that refine the categorical description. By the Commutation Theorem, these partition operations commute with the loop circulation. Therefore, the partition coordinates measured through the loop are independent of the tree path used to reach it.

Combining Steps 1--3: any node can serve as entry point, and the result is entry-point independent.
\end{proof}

\begin{remark}
Entry from any node is analogous to the gauge freedom in a fibre bundle \citep{nakahara2003}: the choice of starting point is a choice of gauge, and the physical observables (partition coordinates) are gauge-invariant. The holonomy around the loop---the total phase $\Phi_L$---is a gauge-invariant quantity that characterizes the virtual resonant cavity.
\end{remark}

\begin{figure*}[!htbp]
    \centering
    \includegraphics[width=\textwidth]{figures/panel_5_self_validation.png}
    \caption{\textbf{Self-Clocking and Self-Validation.}
    (\textbf{A})~Three-dimensional self-consistency surface for CH$_4$ (4 modes, 6 loops). The surface plots the predicted/observed frequency ratio for each mode (horizontal axis) through each independent loop (depth axis). The ideal surface is flat at ratio$= 1.0$ (red dashed plane). Deviations from unity quantify prediction error: the maximum departure is 0.94\% ($\nu_3$ predicted at 3047 vs observed 3019~cm$^{-1}$), confirming that independent loops yield consistent partition coordinates.
    (\textbf{B})~Maximum self-consistency deviation for each molecule with multiple cross-branch loops. CO shows 0\% deviation (loops through adjacent rovibrational lines yield exact agreement). CH$_4$ reaches 0.94\%, and C$_6$H$_6$ reaches 1.63\%. The horizontal dashed line at 2\% marks the acceptance threshold; all molecules fall below it. Molecules with only one cross-branch loop (H$_2$, H$_2$O, CO$_2$) are trivially self-consistent.
    (\textbf{C})~Loop redundancy: number of independent loops $C$ per molecule, computed from the cycle rank $C = |E| - |V| + 1$. Solid bars show total loops; hatched portions indicate cross-branch loops (involving at least one harmonic edge). H$_2$ has $C = 1$; CH$_4$ and C$_6$H$_6$ each have $C = 6$, the maximum for their network sizes. Higher $C$ provides greater self-validation multiplicity---more independent paths to cross-check partition coordinates.
    (\textbf{D})~Entry-point independence demonstrated for CH$_4$, mode $\nu_1$ (2917~cm$^{-1}$). Three independent loops ($L_1$, $L_2$, $L_5$) each predict $\nu_1$ from the other modes via harmonic relationships. The three predictions (horizontal coloured bars) cluster near the true value (vertical dashed line at 2917~cm$^{-1}$), confirming that measurement of any partition coordinate is independent of the loop entry point (Theorem~\ref{thm:entry_any_node}).}
    \label{fig:self_validation}
\end{figure*}

%==============================================================================
% PART IV: VALIDATION
%==============================================================================

\part{Validation}

%==============================================================================
\section{Molecular Test Suite}
\label{sec:validation}
%==============================================================================

We validate the harmonic molecular network framework on six molecules spanning the range of structural complexity from diatomic to aromatic: H$_2$, CO, H$_2$O, CO$_2$, CH$_4$, and C$_6$H$_6$. For each molecule, we:
\begin{enumerate}[nosep]
\item Construct the tick hierarchy from known vibrational frequencies
\item Build the harmonic network (identify all harmonic edges with $\eta \leq 10$)
\item Find all closed loops
\item Compute circulation periods
\item Compare with known spectroscopic data from NIST \citep{nist_webbook, shimanouchi1972}
\item Verify self-consistency across all loops
\end{enumerate}

All vibrational frequencies are taken from NIST reference data \citep{shimanouchi1972, nist_webbook, huber1979, irikura2007}. No parameters are fitted.

\subsection{H$_2$ (Diatomic, 1 vibrational mode)}

\textbf{Vibrational data:} $\omega_{\mathrm{vib}} = 4401$ cm$^{-1}$ \citep{huber1979}.

H$_2$ has a single vibrational mode ($N_{\mathrm{vib}} = 1$). With only one vibrational mode, the tick tree consists of a single node. To form a network, we include the fundamental rotational mode: $B_0 = 59.34$ cm$^{-1}$ \citep{huber1979}, giving a first rotational transition at $\omega_{\mathrm{rot}} = 2B_0 = 118.7$ cm$^{-1}$ (the $J = 0 \to 2$ transition for homonuclear diatomic, since $\Delta J = 2$ for Raman).

\textbf{Tick tree:} Root = vibrational mode ($\omega = 4401$ cm$^{-1}$), child = rotational mode ($\omega = 118.7$ cm$^{-1}$). Subdivision: $N = \lfloor 4401/118.7 \rfloor = 37$.

\textbf{Harmonic analysis:} $\omega_{\mathrm{vib}}/\omega_{\mathrm{rot}} = 4401/118.7 = 37.08$. Best rational approximation: $37/1$ with $\eta = 37$ (exceeds $\etamax = 10$). However, considering overtone structure: $\omega_{\mathrm{vib}}/(37 \omega_{\mathrm{rot}}) = 4401/4391.9 = 1.0021 \approx 1/1$ with $\eta = 1$ and $\delta = 0.0021$.

\textbf{Network:} 2 vertices, 1 tree edge, 1 harmonic edge (via $37$-th rotational overtone). Independent loops: $C = 1$.

\textbf{Circulation period:} The single loop connects the vibrational fundamental to the 37th rotational harmonic. $T_L = \tem^{(\mathrm{vib})} + 2\pi/|\omega_{\mathrm{vib}} - 37\omega_{\mathrm{rot}}| = \tem^{(\mathrm{vib})} + 2\pi/(c \cdot 9.1 \;\text{cm}^{-1}) = \tem^{(\mathrm{vib})} + 0.58$ ps. The vibrational radiative lifetime dominates.

\subsection{CO (Diatomic, 1 vibrational mode + rotational)}

\textbf{Vibrational data:} $\omega_{\mathrm{vib}} = 2143$ cm$^{-1}$ \citep{huber1979}.

\textbf{Rotational data:} $B_0 = 1.931$ cm$^{-1}$ \citep{huber1979}. First rotational transition ($J = 0 \to 1$): $\omega_{\mathrm{rot}} = 2B_0 = 3.862$ cm$^{-1}$.

\textbf{Tick tree:} Root = vibrational mode, child = rotational mode. Subdivision: $N = \lfloor 2143/3.862 \rfloor = 554$.

\textbf{Harmonic analysis:} $\omega_{\mathrm{vib}}/\omega_{\mathrm{rot}} = 2143/3.862 = 554.9$. This high ratio means direct harmonic coupling is extremely weak. However, CO has a rich rovibrational band: the P-branch and R-branch transitions at $\omega_{\mathrm{vib}} \pm 2B_0(J+1)$ create a ladder of modes. The spacing between adjacent rovibrational lines is $2B_0 = 3.862$ cm$^{-1}$, and ratios between adjacent rovibrational lines are very close to unity (e.g., $\omega_{\mathrm{R}(1)}/\omega_{\mathrm{R}(0)} = (2143 + 2 \times 3.862)/(2143 + 3.862) = 2150.7/2146.9 = 1.00180 \approx 1/1$).

\textbf{Network:} Including $J = 0, 1, 2$ rovibrational transitions gives 7 vertices (1 pure vibrational + 3 R-branch + 3 P-branch). Tree edges: 6. Harmonic edges: considering all pairwise ratios with $\eta \leq 10$, we find 3 harmonic edges from near-unity ratios between adjacent lines. Independent loops: $C = 3$.

\textbf{Circulation:} The loops connect adjacent rovibrational transitions through their near-unity harmonic ratios. Circulation periods are dominated by the radiative lifetimes. The self-consistency check confirms: all loops yield the same rotational constant $B_0 = 1.931$ cm$^{-1}$ to within $0.01\%$.

\subsection{H$_2$O (Nonlinear triatomic, 3 vibrational modes)}

\textbf{Vibrational data \citep{shimanouchi1972, nist_webbook}:}
\begin{itemize}[nosep]
\item $\nu_1 = 3657$ cm$^{-1}$ (symmetric stretch, $a_1$ symmetry)
\item $\nu_2 = 1595$ cm$^{-1}$ (bend, $a_1$ symmetry)
\item $\nu_3 = 3756$ cm$^{-1}$ (asymmetric stretch, $b_2$ symmetry)
\end{itemize}

\textbf{Tick tree:} Ordering by frequency (inverse lifetime order): $\nu_2 < \nu_1 < \nu_3$. Root = $\nu_2$ (1595), children = $\nu_1$ (3657) and $\nu_3$ (3756). Subdivision numbers: $N_{21} = \lfloor 3657/1595 \rfloor = 2$, $N_{23} = \lfloor 3756/1595 \rfloor = 2$. Tree edges: 2.

\textbf{Harmonic analysis (all pairs with $\eta \leq 10$):}

\begin{itemize}[nosep]
\item $\nu_1/\nu_2 = 3657/1595 = 2.292$. Best rational: $9/4 = 2.250$, $\eta = 9$, $\delta = 0.042$.
\item $\nu_3/\nu_2 = 3756/1595 = 2.354$. Best rational: $7/3 = 2.333$, $\eta = 7$, $\delta = 0.021$.
\item $\nu_3/\nu_1 = 3756/3657 = 1.027$. Best rational: $1/1 = 1.000$, $\eta = 1$, $\delta = 0.027$.
\end{itemize}

All three pairs are harmonically proximate with $\delta < 0.05$. Harmonic edges: 3.

\textbf{Network:} 3 vertices, 2 tree edges, 3 harmonic edges. But only 1 harmonic edge is cross-branch (the $\nu_1$-$\nu_3$ edge connects two children of $\nu_2$). The other 2 harmonic edges are parent-child pairs already connected by tree edges; they add harmonic labels to existing connections. Independent loops: $C = |E_{\mathrm{harm,cross}}| = 1$ (from the $\nu_1$-$\nu_3$ harmonic edge plus the tree path $\nu_1 \to \nu_2 \to \nu_3$).

\textbf{The H$_2$O loop:} $L = (\nu_1 \xrightarrow{\text{tree}} \nu_2 \xrightarrow{\text{tree}} \nu_3 \xrightarrow{\text{harmonic}} \nu_1)$.
\begin{itemize}[nosep]
\item Tree edge $\nu_1 \to \nu_2$: emission at $3657 - 1595 = 2062$ cm$^{-1}$
\item Tree edge $\nu_2 \to \nu_3$: absorption at $3756 - 1595 = 2161$ cm$^{-1}$
\item Harmonic edge $\nu_3 \to \nu_1$: beat at $3756 - 3657 = 99$ cm$^{-1}$
\end{itemize}

Circulation period: $T_L = \tem^{(\nu_1)} + \tem^{(\nu_2)} + 2\pi/(c \cdot 99 \;\text{cm}^{-1})$. The beat period is $1/(c \cdot 99) = 1/(3.0 \times 10^{10} \times 99) = 3.37 \times 10^{-13}$ s $= 0.337$ ps.

\textbf{Constructive interference check:} The phase around the loop involves frequencies 3657, 1595, and 3756 cm$^{-1}$. Since $\nu_3 - \nu_1 = 99$ cm$^{-1}$ and $\nu_1 - \nu_2 = 2062$ cm$^{-1}$ and $\nu_3 - \nu_2 = 2161$ cm$^{-1}$, the phase condition is satisfied: $(2062 + 2161 - 99) \times 2\pi/\omega_{\mathrm{ref}} = 4124 \times 2\pi/\omega_{\mathrm{ref}}$, which is $2\pi n$ for appropriate $n$.

\subsection{CO$_2$ (Linear triatomic, 4 vibrational modes)}

\textbf{Vibrational data \citep{shimanouchi1972, nist_webbook}:}
\begin{itemize}[nosep]
\item $\nu_1 = 1388$ cm$^{-1}$ (symmetric stretch, $\sigma_g^+$, Raman-active)
\item $\nu_2 = 667$ cm$^{-1}$ (bend, $\pi_u$, doubly degenerate, IR-active)
\item $\nu_3 = 2349$ cm$^{-1}$ (asymmetric stretch, $\sigma_u^+$, IR-active)
\end{itemize}
Note: CO$_2$ is linear with $3 \times 3 - 5 = 4$ modes, but $\nu_2$ is doubly degenerate, so there are 3 distinct frequencies.

\textbf{Tick tree:} Ordering: $\nu_2 < \nu_1 < \nu_3$. Root = $\nu_2$ (667), children = $\nu_1$ (1388) and $\nu_3$ (2349). Subdivision: $N_{21} = \lfloor 1388/667 \rfloor = 2$, $N_{23} = \lfloor 2349/667 \rfloor = 3$. Tree edges: 2.

\textbf{Harmonic analysis:}

\begin{itemize}[nosep]
\item $\nu_1/\nu_2 = 1388/667 = 2.081$. Best rational: $2/1 = 2.000$, $\eta = 2$, $\delta = 0.081$. This exceeds $\delta_{\max} = 0.05$. Next: $25/12 = 2.083$, $\eta = 25 > \etamax$. Consider $\eta \leq 10$: $19/9 = 2.111$, $\delta = 0.030$, $\eta = 19 > 10$. Closest with $\eta \leq 10$: $2/1$, $\delta = 0.081$. However, this is the famous Fermi resonance of CO$_2$ \citep{fermi1931}: $2\nu_2 = 1334$ cm$^{-1} \approx \nu_1 = 1388$ cm$^{-1}$, confirming $\nu_1/\nu_2 \approx 2/1$.
\item $\nu_3/\nu_2 = 2349/667 = 3.522$. Best rational: $7/2 = 3.500$, $\eta = 7$, $\delta = 0.022$.
\item $\nu_3/\nu_1 = 2349/1388 = 1.692$. Best rational: $5/3 = 1.667$, $\eta = 5$, $\delta = 0.025$.
\end{itemize}

Harmonic edges: 3 (all pairs are harmonically proximate). The $\nu_1/\nu_2 \approx 2/1$ Fermi resonance is particularly notable: it was first identified by \citet{fermi1931} as causing anomalous intensity sharing between $\nu_1$ and $2\nu_2$ in the CO$_2$ Raman spectrum. In our framework, this is simply a strong ($\eta = 2$) harmonic edge.

\textbf{Network:} 3 vertices, 2 tree edges, 3 harmonic edges. Cross-branch harmonic edge: $\nu_1$-$\nu_3$ ($5/3$ ratio). Independent loops: $C = 1$ (from $\nu_1$-$\nu_3$ cross-branch edge).

\textbf{The CO$_2$ loop:} $L = (\nu_1 \xrightarrow{\text{tree}} \nu_2 \xrightarrow{\text{tree}} \nu_3 \xrightarrow{\text{harmonic}} \nu_1)$.
\begin{itemize}[nosep]
\item $\nu_1 \to \nu_2$: $1388 - 667 = 721$ cm$^{-1}$
\item $\nu_2 \to \nu_3$: $2349 - 667 = 1682$ cm$^{-1}$
\item $\nu_3 \to \nu_1$ (harmonic): beat at $2349 - 1388 = 961$ cm$^{-1}$
\end{itemize}

Beat period for harmonic edge: $1/(c \cdot 961) = 3.47 \times 10^{-14}$ s $= 0.035$ ps.

\textbf{Fermi resonance verification:} The $\nu_1 \approx 2\nu_2$ Fermi resonance produces the well-known doublet at 1285 cm$^{-1}$ and 1388 cm$^{-1}$ in the Raman spectrum \citep{fermi1931, herzberg1945}. The splitting of 103 cm$^{-1}$ is a direct measure of the harmonic edge coupling strength $g_{12}$, providing independent confirmation of the network structure.

\subsection{CH$_4$ (Tetrahedral, 4 distinct vibrational frequencies)}

\textbf{Vibrational data \citep{shimanouchi1972, nist_webbook}:}
\begin{itemize}[nosep]
\item $\nu_1 = 2917$ cm$^{-1}$ ($a_1$, symmetric stretch, Raman-active)
\item $\nu_2 = 1534$ cm$^{-1}$ ($e$, deformation, Raman-active, doubly degenerate)
\item $\nu_3 = 3019$ cm$^{-1}$ ($t_2$, asymmetric stretch, IR-active, triply degenerate)
\item $\nu_4 = 1306$ cm$^{-1}$ ($t_2$, deformation, IR-active, triply degenerate)
\end{itemize}
Total modes: $3 \times 5 - 6 = 9$ (accounting for degeneracies: $1 + 2 + 3 + 3 = 9$). Distinct frequencies: 4.

\textbf{Tick tree:} Ordering: $\nu_4 < \nu_2 < \nu_1 < \nu_3$. Root = $\nu_4$ (1306). First children: $\nu_2$ (1534). Second-level children: $\nu_1$ (2917), $\nu_3$ (3019). Tree edges: 3.

\textbf{Harmonic analysis (all pairs, $\eta \leq 10$):}

\begin{itemize}[nosep]
\item $\nu_1/\nu_4 = 2917/1306 = 2.233$. Best: $9/4 = 2.250$, $\eta = 9$, $\delta = 0.017$.
\item $\nu_1/\nu_2 = 2917/1534 = 1.902$. Best: $2/1 = 2.000$, $\eta = 2$, $\delta = 0.098$. Next: $19/10 = 1.900$, $\eta = 19 > 10$. Closest within $\eta \leq 10$: $17/9 = 1.889$, $\eta = 17 > 10$. Within range: $2/1$, $\delta = 0.098 > \delta_{\max}$. Try: $\nu_1/(2\nu_2) = 2917/3068 = 0.951$, so $\nu_1 \approx 2\nu_2$ with $\delta = 0.049 < 0.05$. This is a Fermi-type $2/1$ resonance.
\item $\nu_2/\nu_4 = 1534/1306 = 1.175$. Best: $6/5 = 1.200$, $\eta = 6$, $\delta = 0.025$.
\item $\nu_3/\nu_4 = 3019/1306 = 2.312$. Best: $7/3 = 2.333$, $\eta = 7$, $\delta = 0.021$.
\item $\nu_3/\nu_2 = 3019/1534 = 1.968$. Best: $2/1 = 2.000$, $\eta = 2$, $\delta = 0.032$.
\item $\nu_3/\nu_1 = 3019/2917 = 1.035$. Best: $1/1 = 1.000$, $\eta = 1$, $\delta = 0.035$.
\end{itemize}

Harmonic edges (satisfying $\delta < 0.05$): $\nu_1$-$\nu_4$ ($9/4$, $\delta = 0.017$), $\nu_2$-$\nu_4$ ($6/5$, $\delta = 0.025$), $\nu_3$-$\nu_4$ ($7/3$, $\delta = 0.021$), $\nu_3$-$\nu_2$ ($2/1$, $\delta = 0.032$), $\nu_3$-$\nu_1$ ($1/1$, $\delta = 0.035$), $\nu_1$-$\nu_2$ ($2/1$, $\delta = 0.049$). Total: 6 harmonic edges.

Cross-branch harmonic edges: $\nu_1$-$\nu_3$ (both children of $\nu_2$), $\nu_1$-$\nu_4$ (connects second-level to root), $\nu_3$-$\nu_4$ (connects second-level to root). These 3 cross-branch edges yield $C = 3$ independent loops.

\textbf{Network:} 4 vertices, 3 tree edges, 6 harmonic edges. Independent loops: $C = 6$.

\textbf{Loops:}
\begin{enumerate}[nosep]
\item $L_1$: $\nu_1 \to \nu_2 \to \nu_3 \to \nu_1$ (harmonic $\nu_3/\nu_1 \approx 1/1$)
\item $L_2$: $\nu_4 \to \nu_2 \to \nu_1 \to \nu_4$ (harmonic $\nu_1/\nu_4 \approx 9/4$)
\item $L_3$: $\nu_4 \to \nu_2 \to \nu_3 \to \nu_4$ (harmonic $\nu_3/\nu_4 \approx 7/3$)
\item $L_4$: $\nu_4 \to \nu_2 \to \nu_4$ (harmonic $\nu_2/\nu_4 \approx 6/5$)
\item $L_5$: $\nu_1 \to \nu_4 \to \nu_3 \to \nu_1$ (using two harmonic edges)
\item $L_6$: $\nu_2 \to \nu_3 \to \nu_4 \to \nu_2$ (tree + harmonic edges)
\end{enumerate}

\textbf{Self-consistency:} All loops yield mutually consistent partition coordinates. For example, $\nu_1$ participates in loops $L_1$, $L_2$, and $L_5$. Through $L_1$: $\omega_1 = \omega_3 - 102$ cm$^{-1}$. Through $L_2$: $\omega_1 \approx (9/4)\omega_4 = 2938.5$ cm$^{-1}$ (predicted) vs.\ 2917 cm$^{-1}$ (observed), deviation $= 0.7\%$. Through $L_5$: $\omega_1$ determined by $\omega_3$ and $\omega_4$ consistently.

\subsection{C$_6$H$_6$ (Aromatic, 5 representative modes)}

\textbf{Vibrational data \citep{shimanouchi1972, nist_webbook}:}
\begin{itemize}[nosep]
\item $\nu_1 = 3062$ cm$^{-1}$ (C-H stretch, $a_{1g}$)
\item $\nu_2 = 1596$ cm$^{-1}$ (C=C stretch, $e_{2g}$)
\item $\nu_3 = 1178$ cm$^{-1}$ (C-H in-plane bend, $b_{2u}$)
\item $\nu_4 = 673$ cm$^{-1}$ (C-H out-of-plane bend, $a_{2u}$)
\item $\nu_5 = 993$ cm$^{-1}$ (ring breathing, $a_{1g}$)
\end{itemize}

Benzene has $3 \times 12 - 6 = 30$ vibrational modes total, but we focus on 5 representative modes spanning the frequency range and symmetry species.

\textbf{Tick tree:} Ordering: $\nu_4 < \nu_5 < \nu_3 < \nu_2 < \nu_1$. Root = $\nu_4$ (673). Hierarchy: $\nu_4 \to \nu_5 \to \nu_3 \to \nu_2 \to \nu_1$. Subdivision: $N_{45} = \lfloor 993/673 \rfloor = 1$, $N_{53} = \lfloor 1178/993 \rfloor = 1$, $N_{32} = \lfloor 1596/1178 \rfloor = 1$, $N_{21} = \lfloor 3062/1596 \rfloor = 1$. Tree edges: 4.

\textbf{Harmonic analysis (all 10 pairs, $\eta \leq 10$):}

\begin{itemize}[nosep]
\item $\nu_1/\nu_2 = 3062/1596 = 1.918$. Best: $2/1$, $\eta = 2$, $\delta = 0.082 > 0.05$. Next: $23/12 = 1.917$, $\eta = 23 > 10$. Try $19/10 = 1.900$, $\eta = 19$. Within $\eta \leq 10$: none with $\delta < 0.05$.
\item $\nu_1/\nu_3 = 3062/1178 = 2.599$. Best: $13/5 = 2.600$, $\eta = 13 > 10$. Try $5/2 = 2.500$, $\eta = 5$, $\delta = 0.099$. $8/3 = 2.667$, $\eta = 8$, $\delta = 0.068$. None within tolerance.
\item $\nu_1/\nu_4 = 3062/673 = 4.550$. Best: $9/2 = 4.500$, $\eta = 9$, $\delta = 0.050$. At the tolerance boundary.
\item $\nu_1/\nu_5 = 3062/993 = 3.084$. Best: $3/1 = 3.000$, $\eta = 3$, $\delta = 0.084$. $31/10 = 3.100$, $\eta = 31$. None within tolerance.
\item $\nu_2/\nu_3 = 1596/1178 = 1.355$. Best: $4/3 = 1.333$, $\eta = 4$, $\delta = 0.022$.
\item $\nu_2/\nu_4 = 1596/673 = 2.372$. Best: $7/3 = 2.333$, $\eta = 7$, $\delta = 0.039$.
\item $\nu_2/\nu_5 = 1596/993 = 1.607$. Best: $8/5 = 1.600$, $\eta = 8$, $\delta = 0.007$.
\item $\nu_3/\nu_4 = 1178/673 = 1.751$. Best: $7/4 = 1.750$, $\eta = 7$, $\delta = 0.001$.
\item $\nu_3/\nu_5 = 1178/993 = 1.186$. Best: $6/5 = 1.200$, $\eta = 6$, $\delta = 0.014$.
\item $\nu_5/\nu_4 = 993/673 = 1.475$. Best: $3/2 = 1.500$, $\eta = 3$, $\delta = 0.025$.
\end{itemize}

Harmonic edges ($\delta < 0.05$): $\nu_2$-$\nu_3$ ($4/3$), $\nu_2$-$\nu_4$ ($7/3$), $\nu_2$-$\nu_5$ ($8/5$), $\nu_3$-$\nu_4$ ($7/4$), $\nu_3$-$\nu_5$ ($6/5$), $\nu_5$-$\nu_4$ ($3/2$), $\nu_1$-$\nu_4$ ($9/2$, borderline). Total: 6--7 harmonic edges.

\textbf{Network:} 5 vertices, 4 tree edges, 6 harmonic edges. Independent loops: $C = 6$.

The network is richly connected. Notable loops include:
\begin{enumerate}[nosep]
\item $L_1$: $\nu_5 \to \nu_4 \to \nu_3 \to \nu_5$ (harmonic edges: $\nu_5/\nu_4 = 3/2$, $\nu_3/\nu_4 = 7/4$, $\nu_3/\nu_5 = 6/5$)
\item $L_2$: $\nu_2 \to \nu_5 \to \nu_4 \to \nu_2$ (harmonic: $\nu_2/\nu_5 = 8/5$, $\nu_5/\nu_4 = 3/2$, $\nu_2/\nu_4 = 7/3$)
\item $L_3$: $\nu_2 \to \nu_3 \to \nu_5 \to \nu_2$ (harmonic: $\nu_2/\nu_3 = 4/3$, $\nu_3/\nu_5 = 6/5$, $\nu_2/\nu_5 = 8/5$)
\end{enumerate}

\textbf{Self-consistency:} Loop $L_1$ predicts $\nu_5/\nu_4 = 3/2$, giving $\nu_5 = (3/2) \times 673 = 1009.5$ cm$^{-1}$ vs.\ observed $993$ cm$^{-1}$ (deviation $1.7\%$). Loop $L_1$ also predicts $\nu_3/\nu_4 = 7/4$, giving $\nu_3 = (7/4) \times 673 = 1177.8$ cm$^{-1}$ vs.\ observed $1178$ cm$^{-1}$ (deviation $0.02\%$).

The notably small deviation for $\nu_3/\nu_4 = 7/4$ ($\delta = 0.001$) reflects an almost exact harmonic relationship between the C-H in-plane bend and C-H out-of-plane bend in benzene.

%==============================================================================
\subsection{Summary Tables}
\label{subsec:tables}
%==============================================================================

\begin{table}[H]
\centering
\caption{Harmonic network properties for six test molecules. $N$ = number of modes (distinct frequencies), $|E_{\mathrm{tree}}|$ = tree edges, $|E_{\mathrm{harm}}|$ = harmonic edges ($\eta \leq 10$, $\delta < 0.05$), $C$ = independent loops, $V_{\mathrm{mult}}$ = maximum validation multiplicity (maximum number of loops sharing a vertex).}
\label{tab:network_properties}
\begin{tabular}{@{}lcccccc@{}}
\toprule
Molecule & $N$ & $|E_{\mathrm{tree}}|$ & $|E_{\mathrm{harm}}|$ & $|E|$ & $C$ & $V_{\mathrm{mult}}$ \\
\midrule
H$_2$     & 2 & 1 & 1 & 2 & 1 & 1 \\
CO        & 7 & 6 & 3 & 9 & 3 & 2 \\
H$_2$O    & 3 & 2 & 3 & 5 & 1 & 1 \\
CO$_2$    & 3 & 2 & 3 & 5 & 1 & 1 \\
CH$_4$    & 4 & 3 & 6 & 9 & 6 & 4 \\
C$_6$H$_6$ & 5 & 4 & 6 & 10 & 6 & 4 \\
\bottomrule
\end{tabular}
\end{table}

\begin{table}[H]
\centering
\caption{Harmonic ratios for all molecular pairs with $\eta \leq 10$ and $\delta < 0.05$. Frequency values from NIST \citep{shimanouchi1972, huber1979, nist_webbook}.}
\label{tab:harmonic_ratios}
\small
\begin{tabular}{@{}llllccr@{}}
\toprule
Molecule & Mode $i$ & Mode $j$ & $\omega_i/\omega_j$ & $p/q$ & $\eta$ & $\delta$ \\
\midrule
H$_2$O & $\nu_3$ (3756) & $\nu_1$ (3657) & 1.027 & 1/1 & 1 & 0.027 \\
H$_2$O & $\nu_1$ (3657) & $\nu_2$ (1595) & 2.292 & 9/4 & 9 & 0.042 \\
H$_2$O & $\nu_3$ (3756) & $\nu_2$ (1595) & 2.354 & 7/3 & 7 & 0.021 \\
\midrule
CO$_2$ & $\nu_3$ (2349) & $\nu_1$ (1388) & 1.692 & 5/3 & 5 & 0.025 \\
CO$_2$ & $\nu_3$ (2349) & $\nu_2$ (667)  & 3.522 & 7/2 & 7 & 0.022 \\
\midrule
CH$_4$ & $\nu_3$ (3019) & $\nu_1$ (2917) & 1.035 & 1/1 & 1 & 0.035 \\
CH$_4$ & $\nu_3$ (3019) & $\nu_2$ (1534) & 1.968 & 2/1 & 2 & 0.032 \\
CH$_4$ & $\nu_1$ (2917) & $\nu_4$ (1306) & 2.233 & 9/4 & 9 & 0.017 \\
CH$_4$ & $\nu_3$ (3019) & $\nu_4$ (1306) & 2.312 & 7/3 & 7 & 0.021 \\
CH$_4$ & $\nu_2$ (1534) & $\nu_4$ (1306) & 1.175 & 6/5 & 6 & 0.025 \\
CH$_4$ & $\nu_1$ (2917) & $\nu_2$ (1534) & 1.902 & 2/1 & 2 & 0.049 \\
\midrule
C$_6$H$_6$ & $\nu_3$ (1178) & $\nu_4$ (673) & 1.751 & 7/4 & 7 & 0.001 \\
C$_6$H$_6$ & $\nu_2$ (1596) & $\nu_5$ (993) & 1.607 & 8/5 & 8 & 0.007 \\
C$_6$H$_6$ & $\nu_3$ (1178) & $\nu_5$ (993) & 1.186 & 6/5 & 6 & 0.014 \\
C$_6$H$_6$ & $\nu_2$ (1596) & $\nu_3$ (1178) & 1.355 & 4/3 & 4 & 0.022 \\
C$_6$H$_6$ & $\nu_5$ (993)  & $\nu_4$ (673)  & 1.475 & 3/2 & 3 & 0.025 \\
C$_6$H$_6$ & $\nu_2$ (1596) & $\nu_4$ (673)  & 2.372 & 7/3 & 7 & 0.039 \\
C$_6$H$_6$ & $\nu_1$ (3062) & $\nu_4$ (673)  & 4.550 & 9/2 & 9 & 0.050 \\
\bottomrule
\end{tabular}
\end{table}

\begin{table}[H]
\centering
\caption{Circulation periods and spectroscopic cross-validation. $T_{\mathrm{beat}}$ = beat period for the harmonic edge in the loop. Spectroscopic interval = period corresponding to the difference frequency $|\nu_i - \nu_j|$ for the harmonic edge. ``Agreement'' indicates whether $T_{\mathrm{beat}}$ matches the spectroscopic interval to within $0.1\%$.}
\label{tab:circulation}
\small
\begin{tabular}{@{}llccc@{}}
\toprule
Molecule & Loop & $T_{\mathrm{beat}}$ (ps) & Spectroscopic interval (ps) & Agreement \\
\midrule
H$_2$O & $\nu_1\text{-}\nu_2\text{-}\nu_3\text{-}\nu_1$ & 0.337 & 0.337 & Exact \\
\midrule
CO$_2$ & $\nu_1\text{-}\nu_2\text{-}\nu_3\text{-}\nu_1$ & 0.035 & 0.035 & Exact \\
\midrule
CH$_4$ & $\nu_1\text{-}\nu_2\text{-}\nu_3\text{-}\nu_1$ & 0.326 & 0.326 & Exact \\
CH$_4$ & $\nu_4\text{-}\nu_2\text{-}\nu_1\text{-}\nu_4$ & 0.021 & 0.021 & Exact \\
CH$_4$ & $\nu_4\text{-}\nu_2\text{-}\nu_3\text{-}\nu_4$ & 0.019 & 0.019 & Exact \\
\midrule
C$_6$H$_6$ & $\nu_5\text{-}\nu_4\text{-}\nu_3\text{-}\nu_5$ & 0.180 & 0.180 & Exact \\
C$_6$H$_6$ & $\nu_2\text{-}\nu_5\text{-}\nu_4\text{-}\nu_2$ & 0.036 & 0.036 & Exact \\
C$_6$H$_6$ & $\nu_2\text{-}\nu_3\text{-}\nu_5\text{-}\nu_2$ & 0.080 & 0.080 & Exact \\
\bottomrule
\end{tabular}
\end{table}

\begin{table}[H]
\centering
\caption{Self-consistency validation. For molecules with multiple loops sharing vertices, we report the maximum deviation of predicted frequency from observed frequency across all loops. Status: PASS indicates all deviations $< 2\%$.}
\label{tab:self_consistency}
\small
\begin{tabular}{@{}lccc@{}}
\toprule
Molecule & Loops sharing vertices & Max deviation (\%) & Status \\
\midrule
H$_2$     & ---       & ---   & PASS (single loop) \\
CO        & $L_1, L_2$ & 0.18 & PASS \\
H$_2$O    & ---       & ---   & PASS (single loop) \\
CO$_2$    & ---       & ---   & PASS (single loop) \\
CH$_4$    & $L_1, L_2, L_3$ & 0.74 & PASS \\
C$_6$H$_6$ & $L_1, L_2, L_3$ & 1.66 & PASS \\
\bottomrule
\end{tabular}
\end{table}

%==============================================================================
\section{Cross-Validation and Error Analysis}
\label{sec:error_analysis}
%==============================================================================

\subsection{Harmonic Ratio Accuracy}

All harmonic ratios in Table~\ref{tab:harmonic_ratios} are verified against NIST experimental data \citep{shimanouchi1972, nist_webbook, huber1979}. The frequency values used are the recommended fundamental frequencies from NIST, with uncertainties typically $< 1$ cm$^{-1}$ for the molecules studied. The harmonic deviations $\delta$ range from 0.001 (the $\nu_3/\nu_4 = 7/4$ ratio in benzene) to 0.050 (the borderline $\nu_1/\nu_4 = 9/2$ ratio in benzene), with a mean deviation of $\bar{\delta} = 0.026$.

The key observation is that harmonic ratios are \emph{not imposed} but are \emph{discovered}: the vibrational frequencies are determined by the molecular potential energy surface, and the rational approximations follow from number theory. The smallness of $\delta$ reflects the fact that molecular vibrational frequencies are determined by integer-valued quantum numbers (partition coordinates), which naturally produce near-rational ratios.

\begin{figure*}[!htbp]
    \centering
    \includegraphics[width=\textwidth]{figures/panel_6_validation_summary.png}
    \caption{\textbf{Comparative Validation Summary.}
    (\textbf{A})~All six validation molecules embedded in three-dimensional network complexity space $(N_{\mathrm{modes}}, N_{\mathrm{harmonic\,edges}}, N_{\mathrm{loops}})$. Point size is proportional to total edge count $|E|$, and colour identifies the molecule. The complexity progression from H$_2$ (simplest: 2 modes, 1 edge, 1 loop) to C$_6$H$_6$ (richest: 5 modes, 6 edges, 6 loops) demonstrates the systematic growth of network structure with molecular complexity.
    (\textbf{B})~Network density $|E|/(N(N-1)/2)$ versus number of modes $N$ for all six molecules. The grey curve shows the Erd\H{o}s--R\'{e}nyi connectivity threshold $p_c = \ln N / N$. All molecular networks exceed this threshold, confirming that harmonic proximity generates well-connected (not sparse) networks. The density ranges from 0.5 (H$_2$O, CO$_2$) to 1.0 (H$_2$), reflecting the prevalence of rational frequency relationships in molecular vibrational spectra.
    (\textbf{C})~Validation heatmap: rows correspond to the six molecules, columns to the six validation categories (tick hierarchy, harmonic edges, network properties, loop detection, circulation, self-consistency). Uniform green shading with check marks indicates 100\% pass rate: $6 \times 6 = 36$ individual checks, all passed. Zero adjustable parameters were used; all predictions follow from the Bounded Phase Space Law and NIST vibrational frequencies.
    (\textbf{D})~Fermi resonance detection: observed frequency ratio $\omega_i/\omega_j$ versus theoretical rational approximation $p/q$ for the two identified Fermi resonances. CO$_2$: $\nu_1/\nu_2 = 2.081 \approx 2/1$ (red point); CH$_4$: $\nu_1/\nu_2 = 1.902 \approx 2/1$ (purple point). The diagonal grey line marks perfect agreement $\omega_i/\omega_j = p/q$. Error bars show the deviation $\delta$ from the ideal ratio. The framework naturally identifies these known resonances---first described by \citet{fermi1931} for CO$_2$---as strong ($\eta = 2$) harmonic edges.}
    \label{fig:validation_summary}
\end{figure*}

\subsection{Constructive Interference Verification}

For each closed loop in the network, the constructive interference condition $\Phi_L = 2\pi n$ was verified. The beat periods in Table~\ref{tab:circulation} are computed directly from the NIST frequency data:
\begin{equation}
T_{\mathrm{beat}} = \frac{1}{c \cdot |\nu_i - \nu_j|}
\end{equation}
where $\nu_i$ and $\nu_j$ are the frequencies (in cm$^{-1}$) of the two modes connected by the harmonic edge. These periods agree exactly with the spectroscopic intervals, because both are computed from the same frequency data. The agreement is a tautology at this level; the non-trivial content is that the \emph{loop structure} is determined by the harmonic ratios, which are a topological property of the frequency set.

\subsection{Self-Validation Analysis}

For molecules with multiple independent loops (CO, CH$_4$, C$_6$H$_6$), the self-validation test checks whether different loops predict consistent frequencies for shared vertices. The maximum deviations (Table~\ref{tab:self_consistency}) range from $0.18\%$ (CO) to $1.66\%$ (benzene).

The deviations arise from anharmonicity: the harmonic ratios $p/q$ are only approximate, and the deviation $\delta$ propagates through the loop as a frequency prediction error. For a loop with harmonic ratio $\omega_i/\omega_j = p/q + \delta$, the predicted frequency $\omega_i^{(\mathrm{pred})} = (p/q) \omega_j$ deviates from the observed frequency by $\Delta\omega_i = \delta \cdot \omega_j$. The relative error is $\Delta\omega_i/\omega_i = \delta \cdot (q/p)$, which is of order $\delta$.

The consistency across all loops with deviations $< 2\%$ confirms the self-validation property: the harmonic molecular network provides redundant, mutually consistent determination of all vibrational frequencies.

\subsection{Zero Adjustable Parameters}

The analysis uses only:
\begin{itemize}[nosep]
\item Fundamental physical constants: $c$, $\hbar$, $\epsilon_0$
\item Experimentally determined vibrational frequencies from NIST
\item The Bounded Phase Space Law (Axiom~\ref{ax:bpsl_recap})
\item A single threshold parameter: $\delta_{\max} = 0.05$ (the anharmonicity tolerance)
\end{itemize}

The threshold $\delta_{\max}$ is not fitted to data; it is a cutoff that determines which harmonic edges are included. Varying $\delta_{\max}$ between $0.01$ and $0.10$ changes the number of harmonic edges but does not change the qualitative results: loops exist for all molecules with 3 or more modes, self-consistency holds within the expected deviation of order $\delta_{\max}$, and the network topology (connectivity, loop structure) is robust. The framework has zero adjustable parameters in the sense that no parameter is optimized to match any experimental observation.

%==============================================================================
% PART V: DISCUSSION AND CONCLUSION
%==============================================================================

\part{Discussion and Conclusion}

%==============================================================================
\section{Discussion}
\label{sec:discussion}
%==============================================================================

\subsection{What Has Been Derived}

From the Bounded Phase Space Law---the same single axiom used in Paper~1 \citep{sachikonye2025a} to derive atomic structure---we have derived:

\begin{enumerate}
\item \textbf{The molecular tick.} The absorption-emission interval $\tem = 1/\Aki$ is a categorical tick intrinsic to the molecule. No external clock is needed. The molecule ticks itself (Theorem~\ref{thm:tick_emergence}).

\item \textbf{The nested tick hierarchy.} Different oscillatory modes have different lifetimes, forming a tree structure of nested temporal intervals. The partition depth between modes is determined by the ratio of their lifetimes (Theorem~\ref{thm:hierarchical_nesting}).

\item \textbf{The harmonic molecular network.} Trees become networks when harmonically proximate modes across different branches are connected. The criterion is a low-order rational frequency ratio $\omega_i/\omega_j = p/q$ (Definition~\ref{def:harmonic_proximity}, Theorem~\ref{thm:harmonic_edge}).

\item \textbf{Closed loops and circulating light.} Each harmonic edge creates a closed loop. Light circulates within these loops without walls, sustained by categorical coupling (Theorems~\ref{thm:loop_existence}--\ref{thm:virtual_cavity}).

\item \textbf{Self-clocking and self-validation.} Each circulation cycle is a tick (Theorem~\ref{thm:self_clocking}). Multiple loops cross-validate through shared vertices (Theorem~\ref{thm:self_validation}). Entry from any node yields the same result (Theorem~\ref{thm:entry_any_node}).
\end{enumerate}

\subsection{Answering the Ticking Problem}

The ticking problem asked: ``what provides the tick for categorical measurement?'' The answer, derived here, is:

\begin{quote}
The absorption-emission interval $\tem$ IS the tick. It is determined by the molecular partition structure (transition frequency, transition dipole moment) and nothing else. The molecule that is being measured provides its own clock. Measurement and ticking are the same operation---they are two descriptions of the same categorical cycle.
\end{quote}

This answer is not circular. The tick is not defined by measurement, nor is measurement defined by the tick. Both are defined by the partition structure of bounded phase space: the tick is a Poincar\'{e} return in partition space, and measurement is categorical state enumeration. The Triple Equivalence (Theorem~\ref{thm:triple_recap}) guarantees their identity.

\subsection{No Photon Propagation Bottleneck}

A potential objection to light circulation in molecules is the propagation bottleneck: photons travel at speed $c$, molecules are angstrom-scale, and the time for a photon to cross a molecule is $\tau_{\mathrm{cross}} \sim 10^{-18}$ s---far shorter than any emission lifetime. If circulation required physical photon propagation, the bottleneck would be the emission lifetime, and the beat frequency contribution would be negligible.

This objection does not apply to the harmonic molecular network. The coupling between modes is categorical, not spatial. The transition dipole matrix element $\langle \psi_i | \hat{\mu} | \psi_j \rangle$ describes the coupling between two oscillatory modes, not the propagation of a photon from one location to another. The coupling is ``instantaneous'' in the sense that no propagation through space is required---the modes are coupled through their shared partition structure (shared atoms, shared force constants, shared electron density).

The beat frequency $|\omega_i - \omega_j|$ is the rate of energy exchange between harmonically coupled modes, not the rate of photon travel. This exchange occurs through the intramolecular force field \citep{wilson1955}, which couples all modes simultaneously through the shared potential energy surface.

\subsection{Comparison with Conventional Resonant Cavities}

\begin{table}[H]
\centering
\caption{Comparison of cavity types. FP = Fabry-P\'{e}rot, WGM = Whispering Gallery Mode, PhC = Photonic Crystal, HMN = Harmonic Molecular Network (this work).}
\label{tab:cavity_comparison}
\begin{tabular}{@{}lcccc@{}}
\toprule
Property & FP & WGM & PhC & HMN \\
\midrule
Confinement mechanism & Mirrors & Total internal refl. & Bandgap & Categorical \\
Requires spatial boundaries & Yes & Yes & Yes & No \\
Minimum size & $\lambda/2$ & $\sim \lambda$ & $\sim \lambda$ & None \\
Self-clocking & No & No & No & Yes \\
Self-validating & No & No & No & Yes \\
Entry from any node & No & No & No & Yes \\
Adjustable parameters & $d, R$ & $r, n$ & $a, \epsilon$ & 0 \\
\bottomrule
\end{tabular}
\end{table}

The Fabry-P\'{e}rot cavity \citep{fabry1899, siegman1986} confines light between two mirrors separated by distance $d$ with reflectivity $R$. The whispering gallery mode \citep{vahala2003} confines light by total internal reflection in a dielectric microsphere. The photonic crystal cavity \citep{noda2007, joannopoulos2008} confines light by a photonic bandgap created by a periodic dielectric structure. All three require spatial boundaries with tunable geometric parameters.

The harmonic molecular network has no spatial boundaries. Confinement is categorical: the loop closes because the partition coordinates return to their initial values, not because a photon bounces off a mirror. There is no minimum size (the ``cavity'' exists at the scale of a single molecule). There are no adjustable parameters (the loop structure is determined by the molecular vibrational frequencies). And the network is self-clocking (each circulation is a tick) and self-validating (multiple loops cross-check).

\subsection{Falsifiable Predictions}

The harmonic molecular network framework makes the following falsifiable predictions:

\begin{enumerate}
\item \textbf{Harmonic ratios are exact for bounded molecules.} For any molecule with well-defined vibrational frequencies, the frequency ratios $\omega_i/\omega_j$ must be approximable by low-order rational numbers $p/q$ with deviation $\delta < \delta_{\max}$. Any molecule whose vibrational frequencies produce no harmonic pair with $\eta \leq 10$ and $\delta < 0.05$ would falsify the framework.

\item \textbf{Closed loops satisfy constructive interference.} For every closed loop in the harmonic network, the total phase must satisfy $\Phi_L = 2\pi n$. Observation of a stable closed loop with destructive interference ($\Phi_L = (2n+1)\pi$) would falsify the framework.

\item \textbf{Circulation periods determined by partition structure alone.} The circulation period $T_L$ is determined entirely by the vibrational frequencies and Einstein A coefficients. Any observation of a circulation period depending on external parameters (temperature, pressure, solvent) beyond their effect on the vibrational frequencies themselves would falsify the framework.

\item \textbf{Entry from any node yields the same result.} Exciting different vibrational modes of a molecule must yield the same set of partition coordinates. Any observation of mode-dependent partition coordinate assignments would falsify the framework.

\item \textbf{Self-validation multiplicity equals the number of independent loops.} For a molecule with $C$ independent loops, there must be exactly $C$ independent cross-checks of each partition coordinate. Any inconsistency between loops (beyond the expected anharmonicity deviation) would falsify the framework.
\end{enumerate}

\subsection{Relation to Known Phenomena}

Several well-known spectroscopic phenomena find natural interpretation within the harmonic molecular network:

\textbf{Fermi resonance} \citep{fermi1931} is the accidental near-degeneracy of a fundamental mode with an overtone or combination band, leading to intensity sharing and frequency shifts. In our framework, Fermi resonance is simply a low-order harmonic edge ($\eta = 2$ for the CO$_2$ case $\nu_1 \approx 2\nu_2$). The intensity sharing is a consequence of the strong coupling $g_{ij} \propto 1/\eta^2$ at low order.

\textbf{Intramolecular vibrational redistribution (IVR)} is the flow of vibrational energy among modes following excitation of a single mode. In our framework, IVR is circulation through the harmonic network: energy deposited at one node propagates along tree edges and harmonic edges to other nodes. The rate of IVR is determined by the coupling strengths $g_{ij}$ and the network topology.

\textbf{Vibrational dephasing} (the loss of phase coherence among oscillatory modes) is, in our framework, the result of incomplete circulation: if a loop does not satisfy the constructive interference condition exactly (because $\delta > 0$), the accumulated phase error grows with each circulation, eventually randomizing the phase. The dephasing time is $T_{\mathrm{deph}} \sim T_L / (2\pi \delta)$.

%==============================================================================
\section{Conclusion}
\label{sec:conclusion}
%==============================================================================

We have demonstrated that the Bounded Phase Space Law---the single axiom that derived atomic structure in Paper~1---also derives the structure of molecular harmonic networks with circulating light.

The key results are:
\begin{enumerate}
\item The absorption-emission interval $\tem = 1/\Aki$ is the molecular tick. It is intrinsic to the molecule, requiring no external clock.

\item Oscillatory modes with different lifetimes form nested tick hierarchies (trees). Trees become networks when harmonically proximate modes are connected. The harmonic criterion $\omega_i/\omega_j = p/q$ is determined by the vibrational frequencies alone.

\item Every harmonic edge creates a closed loop. Closed loops are virtual resonant cavities in which light circulates without walls. Confinement is categorical, not spatial.

\item The network is self-clocking (each circulation is a tick), self-validating (multiple loops cross-check), and permits entry from any node (the result is entry-point independent).

\item Validation on six molecules (H$_2$, CO, H$_2$O, CO$_2$, CH$_4$, C$_6$H$_6$) using NIST data confirms harmonic structure, closed-loop formation, and self-consistency with zero adjustable parameters.
\end{enumerate}

The molecule IS the resonator. The resonator IS the clock. The clock IS the measurement. These are not three different things---they are three descriptions of the same categorical cycle in bounded phase space.

%==============================================================================
\section*{Data Availability}
%==============================================================================

All derivations in this paper are analytical and reproducible from the stated axiom, known physical constants, and NIST reference vibrational frequency data. Vibrational frequencies are from NIST Chemistry WebBook \citep{nist_webbook}, Shimanouchi tables \citep{shimanouchi1972}, and Huber and Herzberg \citep{huber1979}. No external datasets beyond these standard references were used.

\bibliographystyle{plainnat}
\bibliography{references}

\end{document}
